\documentclass{oupau}
\usepackage{hyperref}

% Comment out for submission
%\usepackage{pdfsync}

%\usepackage{tikz}
%\usetikzlibrary{matrix,arrows} 
%\documentclass[times]{oupau}
%\documentclass[times,doublespace]{oupau}
%\textwidth=0.8\textwidth
%\usepackage{times}

%\usepackage{natbib}
\usepackage{url}
\usepackage[all]{xy}
\bibliographystyle{amsalpha}

\DeclareMathOperator{\alg}{alg}
\DeclareMathOperator{\Aut}{Aut}
\DeclareMathOperator{\Cl}{Cl}
\DeclareMathOperator{\disc}{disc}
\DeclareMathOperator{\End}{End}
\DeclareMathOperator{\Frob}{Frob}
\DeclareMathOperator{\Gal}{Gal}
\DeclareMathOperator{\GL}{GL}
\DeclareMathOperator{\HH}{H}
\DeclareMathOperator{\Hom}{Hom}
\DeclareMathOperator{\Ker}{Ker}
\DeclareMathOperator{\Mat}{Mat}
\DeclareMathOperator{\new}{new}
\DeclareMathOperator{\ord}{ord}
\DeclareMathOperator{\PSL}{PSL}
\DeclareMathOperator{\rank}{rank}
\DeclareMathOperator{\Reg}{Reg}
\DeclareMathOperator{\res}{res}
\DeclareMathOperator{\sat}{sat}
\DeclareMathOperator{\Sel}{Sel}
\DeclareMathOperator{\Tate}{Tate}
\DeclareMathOperator{\Tr}{Tr}
\DeclareMathOperator{\tr}{tr}
\DeclareMathOperator{\tor}{tor}
\DeclareMathOperator{\Vol}{Vol}

\newcommand{\an}{{\rm an}}
\newcommand{\rhobar}{\overline{\rho}}
\newcommand{\h}{\mathfrak{h}}
\newcommand{\hh}{\hat{h}}
\renewcommand{\H}{\HH}


\newcommand{\isom}{\cong}
\newcommand{\tensor}{\otimes}
\newcommand{\vphi}{\varphi}

\newcommand{\CC}{\mathbb{C}}
\newcommand{\cA}{\mathcal{A}}
\newcommand{\cN}{\mathcal{N}}
\newcommand{\FF}{\mathbb{F}}
\newcommand{\OO}{\mathcal{O}}
\newcommand{\QQ}{\mathbb{Q}}
\newcommand{\RR}{\mathbb{R}}
\newcommand{\TT}{\mathbb{T}}
\newcommand{\ZZ}{\mathbb{Z}}
\newcommand{\QQbar}{\overline{\QQ}}

\newcommand{\ncisom}{\approx}   % noncanonical isomorphism

% ---- SHA ----
\DeclareFontEncoding{OT2}{}{} % to enable usage of cyrillic fonts
  \newcommand{\textcyr}[1]{%
    {\fontencoding{OT2}\fontfamily{wncyr}\fontseries{m}\fontshape{n}%
     \selectfont #1}}
\newcommand{\Sha}{{\mbox{\textcyr{Sh}}}}

\newtheorem{conjecture}[theorem]{Conjecture}
\newtheorem{remark}[theorem]{Remark}


\begin{document}
\runningheads{W.A. Stein}{Toward a Generalization of the Gross-Zagier Conjecture}

\title{Toward a Generalization of the Gross-Zagier Conjecture} \author{William
  Stein\affil{1}\footnote{Supported by NSF grant 0555776.}}
\address{\affilnum{1}{Department of Mathematics, University of
    Washington, Web: \url{http://wstein.org}, Email: {\tt wstein@uw.edu}}}


\begin{abstract}
 
  We review some of Kolyvagin's results and conjectures about elliptic
  curves, then make a new conjecture that slightly refines Kolyvagin's
  conjectures.  We introduce a definition of finite index subgroups
  $W_p \subset E(K)$, one for each prime $p$ that is inert in a fixed
  imaginary quadratic field $K$.  These subgroups generalize the group
  $\ZZ y_K$ generated by the Heegner point $y_K \in E(K)$ in the case
  $r_{\an} = 1$.  For any curve with $r_{\an}\geq 1$, we give a
  description of $W_p$, which is conditional on truth of the Birch and
  Swinnerton-Dyer conjecture and our conjectural refinement of
  Kolyvagin's conjecture.  We then deduce the following conditional
  theorem, up to an explicit finite set of primes: (a) the set of
  indexes $[E(K):W_p]$ is finite, and (b) the subgroups $W_p$ with
  $[E(K):W_p]$ maximal satisfy a higher-rank generalization of the
  Gross-Zagier formula.  We also investigate a higher-rank
  generalization of a conjecture of Gross-Zagier.

\end{abstract}

\keywords{number theory; Birch and Swinnerton-Dyer Conjecture; Euler
  systems; Heegner points; Kolyvagin's conjecture; computational
  number theory; Gross-Zagier theorem}

\received{}
\maketitle

\section{Introduction}\label{sec:intro}
Let $E$ be an elliptic curve defined over $\QQ$.  
The order of vanishing $r_{\an}$ at $s=1$ of the Hasse-Weil 
$L$-series $L(E/\QQ,s)$ of $E$ is defined because $E$ is modular 
(see \cite{breuil-conrad-diamond-taylor, wiles:fermat}).
The Birch and
Swinnerton-Dyer (BSD) rank conjecture \cite{birch:edsac} asserts that
$r_{\an}$ is equal to  the algebraic rank $r_{\alg}$ of $E(\QQ)$.
The BSD formula then gives a conjectural formula for the
leading coefficient of the Taylor expansion about $s=1$ of
$L(E/\QQ,s)$; this formula resembles the analytic class number formula.  The
BSD rank conjecture is known for curves with $r_{\an}\leq 1$, but
there has been relatively little progress toward the BSD rank
conjecture when $r_{\an} \geq 2$.

In the late 1980s, Kolyvagin wrote several landmark papers that
combined the Gross-Zagier theorem \cite{gross-zagier} about heights of
Heegner points over quadratic imaginary fields $K$, a theorem
\cite{bump-friedberg-hoffstein:nonvanishing} about nonvanishing of
special values of twists of $L$-functions, and relations involving
Hecke operators between Heegner points over ring class fields of $K$
to prove that if $r_{\an}\leq 1$, then the BSD rank conjecture is true
for $E$.  Kolyvagin wrote \cite{kolyvagin:structure_of_selmer} on the
case of general rank, in which he computes the elementary invariants
of the Selmer groups of any elliptic curve $E$ of any rank in terms of
properties of Heegner points, assuming a certain nontriviality
hypothesis.  It was until recently unclear whether or not this
hypothesis was ever satisfied for any curve with $r_{\an}\geq 2$.
Fortunately, this hypothesis has now been confirmed numerically (with
high probability) in one case of a rank 2 curve
\cite{jetchev-lauter-stein}.  
%{\em The aim of the present paper is to
%  take another small step toward a Kolyvagin-style approach to BSD for
%  curves with $r_{\an}\geq 2$ by studying extra structure on $E$ that
%  we get by assuming conjectures of Birch, Swinnerton-Dyer, and
%  Kolyvagin.}

We review some of Kolyvagin's results and conjectures from
\cite{kolyvagin:structure_of_selmer}, then make a new conjecture that
refines Kolyvagin's conjectures.  Using reduction modulo $p$ of
Heegner points, we introduce a definition of finite index subgroups
$W_p \subset E(K)$, one for each prime $p$ that is inert in $K$.  Let
$y_K\in E(K)$ be the associated Heegner point as in
Equation~\eqref{eqn:yK} below.  Then these subgroups $W_p$ generalize
the group $\ZZ y_K$ in the case $r_{\an} = 1$.  For any $r_{\an}\geq
1$, we give a description of $W_p$, which is conditional on truth of
the BSD conjecture and our conjectural refinement of Kolyvagin's
conjecture.  We then deduce the following conditional theorem (see
Theorems~\ref{thm:refined_evidence} and \ref{thm}), up to an explicit
finite set of primes: (a) the set of indexes $[E(K):W_p]$ is finite,
and (b) the subgroups $W_p$ with $[E(K):W_p]$ maximal 
 satisfy a higher-rank generalization of the
Gross-Zagier formula (see \eqref{eq:ggz} below).  We also give
numerical data and a new conjecture about the existence of
Gross-Zagier subgroups.

We leave open far more questions than we answer, and we intend to
follow up on these questions in subsequent papers.  For example,
perhaps the definition of the groups $W_p$ can be refined and
generalized in various ways, and results similar to those in this
paper proved about them.  It would be interesting to find a practical
algorithm that can provably compute the groups $W_p$ for a particular
$p$, assuming that $E(K)$ has already been computed.  We also hope to
find a higher-rank analogue of the Gross-Zagier formula over the
Hilbert class field of $K$, involving the Petersson inner product,
modular forms, and Rankin-Selberg convolutions $L_{\mathcal{A}}(f,s)$,
as in \cite{gross-zagier}, which is consistent with the results we
prove about the groups $W_p$ in this paper.  It would also be valuable
to give proofs of the results of \cite{kolyvagin:structure_of_selmer}
building on \cite{mccallum:kolyvagin} instead of
\cite{kolyvagin:structureofsha}, possibly using results from the
present paper.

We briefly outline the structure of this paper. In the first few
sections, we state the BSD conjecture and Gross-Zagier formula, define
Kolyvagin points, state Kolyvagin's conjectures, and then define
certain finite index subgroups $W_p$ of $E(K)$.  In the rest of the
paper, we study reduction mod $p$, conditionally deduce the structure
of $W_p$, and give some numerical examples.

More precisely, we do the following.  In Section~\ref{sec:ggz} we
state the full Birch and Swinnerton-Dyer conjecture over an imaginary
quadratic field $K$, and state a generalized Gross-Zagier formula for
elliptic curves of any rank.  In Section~\ref{sec:heegner}, we
introduce the Kolyvagin points $P_{\lambda}$ on $E$ over ring class
fields of $K$, and deduce some key properities of these points.  We
state Kolvagin's conjectures from \cite{kolyvagin:structure_of_selmer}
along with some of their consequences in Section~\ref{sec:kolyconj}.
We also state a conjecture that refines Kolyvagin's conjectures and
also refines a conjecture of Gross-Zagier.  In
Section~\ref{sec:kolysubgroups} we use reductions of Kolyvagin points
to define, for every prime $p$ that is inert in $K$, a finite index
subgroup $W_p$ of $E(K)$.  Section~\ref{sec:reduction} lays some
general foundations for our later determination of the structure of
$W_p$ by studying the image of a fixed $Q\in E(K)$ in
$E(\FF_{p^2})/(p+1)$.  Section~\ref{sec:maxwp} presents a conditional
proof that (up to primes not in the set $B(E)$) maximal index
subgroups exist and that they satisfy our generalized Gross-Zagier
formula.  Finally, in Section~\ref{sec:gzexist} we numerically
investigate the existence of Gross-Zagier subgroups of $E(K)$, and
give evidence for a higher-rank generalization of a conjecture of
Gross-Zagier.


\vspace{2ex}
\noindent{\bf Acknowledgement:} We thank R. Bradshaw, K.
Buzzard, C. Citro, J. Coates, C. Cornut, M. Flach, R. Greenberg,
B. Gross, D. Jetchev, K. Lauter, B. Mazur, R. Miller, and 
Tonghai Yang for helpful conversations.  We thank Amod
Agashe and Andrei Jorza for carefully reading a draft of the paper
and providing many helpful comments, and we thank the anonymous
referee for much helpful feedback.

\begin{quote}
``It is always good to try to prove true theorems.''

-- Bryan Birch
\end{quote}

%\tableofcontents


\subsection{Notation and Conventions}\label{notation}
Let $A$ be an abelian group.  Let $A_{\tor}$ be the subgroup of
elements of $A$ of finite order and let $A_{/\tor} = A/A_{\tor}$ denote the
quotient of $A$ by its torsion subgroup.  Let $A[n]$ be the
subgroup of elements of $A$ of order $n$, and for any prime $\ell$,
let $A(\ell)$ be the subgroup of elements of $\ell$-power order.  For
$z\in A$, let $e=\ord_{\ell}(z)$ be the largest integer $e$ such that $z
= \ell^e y$ for some $y \in A$, or $\ord_{\ell}(z)=\infty$ if the set
of $e$ is unbounded.
If $a_1,\ldots, a_n$ are elements of an additive or multiplicative
group $A$, we let $\langle a_1, \ldots , a_n\rangle$ 
denote the subgroup of $A$ generated by the $a_i$.

Throughout this paper, $E$ denotes an elliptic curve defined over
$\QQ$ of conductor $N$, and $K$ is a quadratic imaginary field with
$D=\disc(K)$ coprime to $N$ that satisfies the {\em Heegner
  hypothesis}---each prime dividing $N$ splits in $K$.  We fix an
ideal $\mathcal{N}$ in $\OO_K$ such that $\OO_K/\mathcal{N}$ is cyclic
of order $N$.  Let $H$ be the Hilbert class field of $K$,
let $\pi:X_0(N)\to E$ be a fixed choice of modular 
parametrization (see Section~\ref{sec:heegner} below),
and let 
\begin{equation}\label{eqn:yK}
  y_K = \Tr_{H/K}(\pi((\CC/\OO_K,\mathcal{N}^{-1}/\OO_K))) \in E(K)
\end{equation}
be the Heegner point associated to $K$.

Let $c$ denote the Manin constant of $E$ (see
Section~\ref{sec:ggz}), and $c_q$ the Tamagawa numbers of $E$ at
primes $q\mid N$.  Unless otherwise stated, everywhere in this paper
$p$ denotes a prime that is inert in $K$.

\section{Gross-Zagier Subgroups}\label{sec:ggz}
In this section, we fix our notation and conventions, and define
the Manin constant.   Then we recall the statement of the full
Birch and Swinnerton-Dyer conjecture over an imaginary quadratic
field $K$.  We give a new definition of {\em Gross-Zagier subgroups}
of $E(K)$ and prove that they all satisfy a Gross-Zagier style
formula.  When $r_{\an}=1$, we prove that $\ZZ y_K$ is the
unique Gross-Zagier subgroup, up to torsion. 

Let $E$ be an elliptic curve over $\QQ$ and let $K$ be a quadratic
imaginary field that satisfies the Heegner hypothesis -- so $K$ has
discriminant $D<-4$, each prime dividing the conductor $N$ of
$E$ splits in $K$, and $\gcd(D,N)=1$.  Let $\OO_K$ be the ring of
integers of $K$.  Let $E^D$ denote the quadratic twist of $E$ by $D$.
Throughout this paper, except briefly in Section~\ref{sec:gzexist}, we
{\em always} assume that
\begin{equation}\label{eq:minrank}
r_{\an}(E/\QQ) > r_{\an}(E^D/\QQ) \leq 1
\end{equation}

Recall that under the Heegner hypothesis the sign of the functional equation of 
$$
 L(E/K,s) = L(E/\QQ,s) \cdot L(E^D/\QQ,s)
$$ is $-1$,
so the sign in the functional equations for 
$L(E/\QQ,s)$ and $L(E^D/\QQ,s)$ are different, hence
$$\ord_{s=1}L(E/\QQ,s) \not\equiv \ord_{s=1}L(E^D/\QQ,s) \pmod{2}.$$

\begin{proposition}
  Suppose $E$ is an elliptic curve with $r_{\an}(E/\QQ)>0$. Then there
  exist infinitely many $D$ satisfying the Heegner hypothesis with
$$ 
r_{\an}(E/\QQ) > r_{\an}(E^D/\QQ) \leq 1.
$$
\end{proposition}
\begin{proof}
  The main theorem of
  \cite{bump-friedberg-hoffstein:nonvanishing}
implies the existence of infinitely many
$D$ with $r_{\an}(E^D/\QQ) \leq 1$. 
Since $r_{\an}(E/\QQ)>0$  and  $r_{\an}(E/\QQ) \not\equiv r_{\an}(E^D/\QQ)\pmod{2}$, 
the inequality  $r_{\an}(E/\QQ) > r_{\an}(E^D/\QQ)$ also holds. 
\end{proof}

Let
$\omega = 2\pi i c f(z) dz$ be the pullback of a minimal invariant
differential on $E$, where $f(z) \in S_2(\Gamma_0(N))$ is a cuspidal
newform, and $c$ is the Manin constant of $E$ (see
\cite{agashe-ribet-stein:manin}).  For each prime $q\mid N$, let $c_q$
be the Tamagawa number of $E$ at $q$.  Set
$r=r_{\an}(E/K)=\ord_{s=1}L(E/K,s)$, which is defined since every
elliptic curve over $\QQ$ is modular.  Let $\|\omega\|^2 =
\int_{E(\CC)}\omega\wedge \overline{i\omega} = 2\cdot
\Vol(\CC/\Lambda)$.  The Shafarevich-Tate group of $E$ over a
number field $M$ is 
$$
 \Sha(E/M) = \ker\left(\H^1(M,E) \to \bigoplus_{v} H^1(M_v, E)\right).
$$

The following is a formulation of the Birch and Swinnerton-Dyer
conjecture \cite[pg.~311]{gross-zagier} over $K$.
\begin{conjecture}[Birch and Swinnerton-Dyer]\label{conj:bsd}
  The Mordell-Weil group $E(K)$ has rank $r=\ord_{s=1} L(E/K,s)$, the
  Shafarevich-Tate group $\Sha(E/K)$ is finite, and
\begin{equation}\label{eq:bsd}
\frac{L^{(r)}(E/K, 1)}{r!}
 = \frac{\#\Sha(E/K) \cdot \|\omega\|^2 \cdot \Reg(E/K) \cdot 
    \left(\prod_{q\mid N}c_q\right)^2}
   {\# E(K)_{\tor}^2 \cdot \sqrt{|D|}}.
\end{equation}
\end{conjecture}

Let $\Sha_\an$ be the order of $\Sha(E/K)$ that is
predicted by Conjecture~\ref{conj:bsd}.  
The existence of the Cassels-Tate pairing implies that if $\Sha(E/K)$
is finite, then $\#\Sha(E/K)$ is a perfect square, so 
Conjecture~\ref{conj:bsd} implies that $\sqrt{\Sha_{\an}}$ is
an integer.   Recall from Section~\ref{notation}
that $A_{/\tor} = A/A_{\tor}$. 
\begin{definition}[Gross-Zagier subgroup]\label{defn:ggz}
A {\em Gross-Zagier subgroup} $W\subset E(K)$ 
is a torsion-free subgroup such that $E^D(\QQ) \subset W + E(K)_{\tor}$, 
the quotient $E(K)_{/\tor} / W$ is cyclic, and
\begin{equation}\label{eq:ggzsubgroup}
 [E(K):W] = c \cdot \prod c_q\cdot \sqrt{\Sha_{\an}}.
\end{equation}
\end{definition}
For any set $S$ of primes, we say that a subgroup $W\subset
E(K)$ is a {\em Gross-Zagier subgroup up to primes not in $S$} if $W$
has no $p$-torsion for $p\not\in S$ and all the conditions
of Definition~\ref{defn:ggz} holds up to primes not in $S$.

We will numerically investigate the existence of Gross-Zagier
subgroups in Section~\ref{sec:gzexist}, assuming that
Conjecture~\ref{conj:bsd} is true.  Even the existence of Gross-Zagier
subgroups of every $E(K)$ is far from clear, since if they exist, then
$\#E(K)_{\tor}$ divides $c\prod c_q \cdot \sqrt{\Sha_{\an}}$.  In
fact, we will give an example of an $E(K)$ that does not have any
Gross-Zagier subgroups (this example does not satisfy
\eqref{eq:minrank}).

In the following proposition we do not assume the
Conjecture~\ref{conj:bsd}.  Thus $\Sha_{\an}$ {\em a priori} could
just be some meaningless transcendental number.  Also, for any
subgroup $H\subset E(K)$, we write $\Reg(H)$ for the absolute value of
the determinant of the height pairing matrix on any basis for $H$
modulo torsion.
\begin{proposition}\label{prop:ggz}
If $W$ is a  Gross-Zagier subgroup, then $W$ satisfies the
  {\em generalized Gross-Zagier formula}:
\begin{equation}\label{eq:ggz}
\frac{L^{(r)}(E/K,1)}{r!}
 = 
        \frac{\|\omega\|^2}{c^2 \cdot  \sqrt{|D|}}
     \cdot \Reg(W).
\end{equation}
More generally, a torsion-free subgroup $W\subset E(K)$ satisfies the
generalized Gross-Zagier formula if and only if it has index $c \cdot
\prod c_q\cdot \sqrt{\Sha_{\an}}$ in $E(K)$.
\end{proposition}
\begin{proof}
  The BSD formula \eqref{eq:bsd} with $\#\Sha(E/K)$ replaced by
  $\Sha_{\an}$ implies that \eqref{eq:ggz} holds if and only if
\begin{equation}\label{eqn:bigregv}
\frac{\|\omega\|^2}{c^2 \cdot  \sqrt{|D|}}
     \cdot \Reg(W)
 = \frac{\Sha_{\an} \cdot \|\omega\|^2 \cdot \Reg(E/K) \cdot  
    \left(\prod_{p\mid N}c_q\right)^2}
   {\# E(K)_{\tor}^2 \cdot \sqrt{|D|}}.
\end{equation}
Our hypotheses that  $[E(K):W]$ is finite and that $W$ is torsion free
imply that
\begin{equation}\label{eqn:bigregv2}
 [E(K) : W]^2 = 
    \frac{\Reg(W) \cdot \#E(K)_{\tor}^2}
         {\Reg(E/K)}
\end{equation}
Manipulate \eqref{eqn:bigregv} by cancelling everything in common on
both sides and putting the regulators and torsion on the left, and
everything else on the right.  The substitution \eqref{eqn:bigregv2}
then shows that $ [E(K):W]^2 = c^2 \cdot \left(\prod_{q\mid
    N}c_q\right)^2 \cdot \Sha_{\an} $ if and only if \eqref{eq:ggz}
holds.  Taking square roots proves the proposition.  \qed\end{proof}

\begin{corollary}\label{cor:gzrank1}
  Let $y_K \in E(K)$ be the Heegner point after fixing a choice of
  ideal $\cN$ as in Equation~\eqref{eqn:yK}, and assume that $E$ has
  analytic rank $1$.  Then the Gross-Zagier subgroups of $E(K)$ are
  the cyclic groups $\langle y_K+P\rangle $, for all $P\in
  E(K)_{\tor}$.
\end{corollary}
\begin{proof}
  By \cite{kolyvagin:subclass}, $E(K)$ is of rank $1$, and by
  Proposition~\ref{prop:ggz} the Gross-Zagier formula \cite[Thm.~2.1,
  pg.~311]{gross-zagier} implies that $[E(K): \langle y_K\rangle] =
  c\prod c_q \sqrt{\Sha_{\an}}$ (see also, \cite[Conj.~2.2,
  pg.~311]{gross-zagier}).  Since $E(K)_{/\tor}$ is free of rank $1$
  and $\langle y_K\rangle $ is torsion free, $E(K)_{/\tor} / \langle
  y_K\rangle$ is cyclic, so $\langle y_K\rangle$ is a Gross-Zagier
  subgroup. The same argument proves this with $y_K$ replaced by
  $y_K+P$ for any $P\in E(K)_{\tor}$, since $y_K$ and $y_K+P$ have the
  same height.  If~$W$ is any Gross-Zagier subgroup, then since $E(K)$
  has rank one we must have $W \equiv \langle y_K\rangle
  \pmod{E(K)_{\tor}}$, so $W = \langle y_K + P \rangle$ for some $P\in
  E(K)_{\tor}$.
\end{proof}

\section{Heegner and Kolyvagin Points}\label{sec:heegner}
In this section, we define certain subsets $\Lambda^k_{\ell^n}\subset
\ZZ$ of positive square-free integers.  For each integer $\lambda \in
\Lambda^k_{\ell^n}$, we consider the corresponding ring class field
$K_{\lambda}$, and we define elements $I_{\lambda}, J_{\lambda} \in
\ZZ[\Gal(K_{\lambda}/K)]$.  We then apply these group ring elements to
the Heegner points $y_{\lambda} \in E(K_{\lambda})$ to obtain the
Kolyvagin points $P_{\lambda} \in E(K_{\lambda})$.  Finally, we prove
the $\Gal(K_{\lambda}/K)$-equivariance of the equivalence class
$P_{\lambda} + \ell^n E(K)$ in $E(K)/\ell^n E(K)$. 

For any integer $m$, let $a_m = a_m(E)$ be the $m$th coefficient
of the $L$-series $\sum a_m/m^s$ attached to $E$.
Let $\ell$ be any prime and $n$ any positive integer.  
For any nonnegative integer $k$, 
let $\Lambda_{\ell^n}^k$ be the
set of squarefree positive integers $\lambda = p_1\dots p_k$ coprime to $N$, 
where each $p_i$ is inert in $K$ and
$$
   a_{p_i} \equiv p_i + 1 \equiv 0 \pmod{{\ell^n}}.
$$
When $k=0$, we set $\Lambda^0_{\ell^n} = \{1\}$.  
The Chebotarev density theorem implies that $\Lambda^k_{\ell^n}$ is infinite 
for any $k\geq 1$.

Recall from Section~\ref{notation} that we fixed an ideal
$\mathcal{N}$ in $\OO_K$ such that $\OO_K/\mathcal{N}$ is cyclic of
order $N$, and let $\OO_{\lambda} = \ZZ + \lambda \OO_K$ be the order
in $\OO_K$ of conductor $\lambda$. 
Let $X_0(N)$ be the compact modular curve defined over $\QQ$ that classifies
isomorphism classes of elliptic curves equipped with a cyclic subgroup
of order $N$. 
Fix a choice of minimal modular parametrization $\pi:X_0(N)\to E$, which
exists by the modularity theorem \cite{breuil-conrad-diamond-taylor, wiles:fermat}.
For each $\lambda \in \Lambda_{\ell^n}^k$, the {\em Heegner point}
$$
  x_{\lambda} = [(\CC/\OO_{\lambda}, (\mathcal{N}\cap\OO_{\lambda})^{-1} / \OO_{\lambda})]
                \in X_0(N)(K_{\lambda})
$$
is defined over the ring class field $K_{\lambda}$ of $K$ of conductor $\lambda$. 
\begin{definition}[Heegner point]
The Heegner point $y_{\lambda}$ associated to $\lambda \in \Lambda_{\ell^n}^k$
is 
$$
  y_{\lambda} = \pi(x_{\lambda}) \in E(K_{\lambda}).
$$
\end{definition}
We emphasize that that $y_{\lambda}$ depends on the choice of modular
parametrization $\pi_E$ and the ideal $\cN$ in $\OO_K$ with
$\OO_K/\cN=\ZZ/N\ZZ$.   However, once we fix that data, the Heegner points
for all $\lambda$ are defined. 

For $\lambda \in \Lambda_{\ell^n}^k$, let $ G_{\lambda} = \Gal(K_{\lambda}/K_1) $
and note that we have a canonical isomorphism
$$
 G_{\lambda} \isom \prod_{p \mid \lambda} G_p,
$$
where the group $G_p = \Gal(K_p/K_1) = \langle t_p \rangle$ is cyclic of order $p+1$,
with some (non-canonical) choice $t_p$ of generator.  Let
$$
I_p = \sum_{i=1}^p i t_p^i \in \ZZ[G_p]
\quad\text{ and }\quad
I_{\lambda} = \prod_{p\mid \lambda} I_p \in \ZZ[G_{\lambda}].
$$
Let $R$ be a set of representatives in $\Gal(K_{\lambda}/K)$
for the quotient group
$\Gal(K_{\lambda}/K) / \Gal(K_{\lambda}/K_1) \isom \Gal(K_1/K)$, 
and  let
$$
J_{\lambda} = \sum_{g \in R} g \in \ZZ[G_{\lambda}].
$$
\begin{definition}[Kolyvagin Point]
The Kolyvagin point $P_{\lambda}$ associated to $\lambda \in \Lambda^k_{\ell_n}$ is
$$
  P_{\lambda} = J_{\lambda} I_{\lambda} y_{\lambda} \in E(K_{\lambda}).
$$
\end{definition}
Note that $P_1 = y_K \in E(K)$. 



Let $R=\End(E/\CC)$ and let $B(E)$\label{defn:be} be the set of {\em odd}
primes $\ell$ that do not divide $\disc(R)$ and such that the
$\ell$-adic representation
$\Gal(\QQbar/\QQ)\to\Aut_{R}(\Tate_{\ell}(E))$ is surjective.  By a
theorem of Serre \cite{serre:propgal}, the set $B(E)$ contains all but
finitely many primes (see \cite{bsdalg1} for algorithms to bound
$B(E)$).  Let $T_p$ be the $p$th Hecke operator on the Jacobian $J_0(N)$
of $X_0(N)$, and for each prime $p\mid \lambda$, let $\Tr_p$ be the
trace $J_0(N)(K_{\lambda}) \to J_0(N)(K_{\lambda/p})$.

\begin{proposition}\label{prop:eulersys}
  The points $y_{\lambda}$ form an Euler system, in the sense that if
  $\lambda = p \lambda'$ for a prime~$p$ and $\lambda \in
  \Lambda_{\ell}$, then $ y_{\lambda} = \Frob_{\wp}(y_{\lambda'})
  \pmod{\wp} $ for all primes $\wp$ of $K_{\lambda}$ over $p$, and
  $\Tr_p(x_{\lambda}) = T_p (x_{\lambda'}) $ in $J_0(N)$.
\end{proposition}
\begin{proof}
See \cite[Prop.~3.7]{gross:kolyvagin}.
\end{proof}


\begin{proposition}\label{prop:invariant}
We have
$$
  [I_{\lambda} y_{\lambda}] \in (E(K_{\lambda}) / {\ell^n} E(K_{\lambda}))^{G_{\lambda}}
\qquad\text{ and }
\qquad 
  [P_{\lambda}] \in (E(K_{\lambda}) / {\ell^n} E(K_{\lambda}))^{\Gal(K_{\lambda}/K)}
$$
\end{proposition}
\begin{proof}
  Though standard (see, e.g., \cite[Prop.~3.6]{gross:kolyvagin}) this
  proposition plays a key role in Section~\ref{sec:kolysubgroups}, so
  we give a proof here for the convenience of the reader.  The first
  statement implies the second, since $[P_{\lambda}]$ is the
  $\Gal(K_{1}/K)$ trace of $[I_{\lambda} y_{\lambda}]$.  It remains to
  prove the first inclusion. For this, it suffices to show that
  $[I_{\lambda} y_{\lambda}]$ is fixed by $t_p$ for all primes $p\mid \lambda$, as
  these elements generate $G_{\lambda}$.  We will prove this by
  showing that $(t_p - 1) I_{\lambda} y_{\lambda}$ lies in ${\ell^n}
  E(K_{\lambda})$.

Write $\lambda = p \lambda'$.  We have
\begin{equation}\label{iprel}
 (t_p - 1) I_p =  (t_p - 1) \cdot \left(\sum_{i=1}^p i t_p^i\right)
               = p + 1 - \Tr_p,
\end{equation}
where as above $\Tr_p = \Tr_{K_{\lambda}/K_{\lambda'}}$.
Note that this is the only place in the proof where we use the explicit
definition of $I_p$ as $\sum_{i=1}^p i t_p^i$, and in fact
we could instead replace $I_p$ by any element $I$ of $\ZZ[G_{\lambda}]$
such that 
$$
(t_p - 1) I = p + 1 - \Tr_p,
$$
but doing so does not seem to lead to anything interesting.
Note that the Euler system relation (see Proposition~\ref{prop:eulersys})
and our hypothesis that $a_p \equiv 0 \pmod{{\ell^n}}$ together
imply that
$$
 \Tr_{p} I_{\lambda'} y_{\lambda} = I_{\lambda'} \Tr_{p} y_{\lambda}= I_{\lambda'} a_p y_{\lambda'} \in {\ell^n} E(K_{\lambda}).
$$ 
We have
$$
  (t_p  - 1) I_{\lambda} = (t_p - 1) I_{p} I_{\lambda'}
                         = (p + 1 - \Tr_p) I_{\lambda'}
$$
in $\ZZ[G_{\lambda}]$, so since $p+1\equiv 0 \pmod{{\ell^n}}$
$$
  (t_p - 1) I_{\lambda} y_{\lambda} 
     = (p + 1) I_{\lambda'} y_{\lambda} - \Tr_p I_{\lambda'} y_{\lambda}
     \in \ell^n E(K_{\lambda}).
$$
\qed\end{proof}

\section{Kolyvagin's Conjectures and their Consequences}\label{sec:kolyconj}
For any prime $\ell$ and positive integer $n$, let 
$$
  \displaystyle \Lambda_{\ell^n} = \bigcup_{\text{all } k \geq 0} \Lambda_{\ell^n}^k
$$ 
be the set of square-free positive integers $\lambda$ such that
$\ell^n\mid \gcd(a_p, p+1)$ for each $p\mid \lambda$.  In this
section, we define maps $n, m: \Lambda_{\ell} \to \ZZ \cup \{\infty\}$
that measure $\ell$-divisibility properties of $\lambda$ and
$P_{\lambda}$ for all $\lambda \in \Lambda_{\ell}$.  We state
Kolyvagin's ``Conjecture A'' that there exists $\lambda$ with
$m(\lambda)\neq \infty$, then state Kolyvagin's structure theorem,
which describes the structure of $\Sel^{(\ell^b)}(E/K)$, for $b$
sufficiently large, in terms of the maps $n$ and $m$.  Finally, we
state Kolyvagin's stronger ``Conjecture D'', which basically asserts
that if $f$ is the smallest nonnegative integer such that
$m(\lambda)\neq \infty$ for some $\lambda \in \Lambda_{\ell}^f$, then
for sufficiently large $k$ the cohomology classes
$\tau_{\lambda,\ell^n}$ with $\lambda \in \Lambda_{\ell^n+k}^f$
generate a subgroup of $\Sel^{(\ell^n)}(E/K)$ that equals the image of
a subgroup $V$ of $E(K)$.  To motivate Conjecture~\ref{conj:D}, we
prove that it implies that $\rank(E(\QQ)) = f+1$ and $\Sha(E/K)(\ell)$
is finite for each $\ell \in B(E)$ and determine the structure of $V$
(see Proposition~\ref{prop:conjD}).

Recall that we defined $\ord_{\ell}$ in Section~\ref{notation}.
Define two set-theoretic maps 
$$
 n, m: \Lambda_{\ell} \to \ZZ \cup \{\infty\}
$$
by
$$
  n(\lambda) = \max\{e : \lambda \in \Lambda_{\ell^e}\}\
\qquad\text{ and }\qquad
  m(\lambda) = \ord_{\ell}([P_{\lambda}]),
$$
where $[P_{\lambda}]$ denotes the equivalence
class of $P_{\lambda}$ in $E(K_{\lambda})/\ell^{n(\lambda)}
E(K_{\lambda})$.  For each integer $k\geq 0$, let 
$$
  m_{\ell,k} = \min(m(\Lambda^k_{\ell}))
\qquad\text{and}\qquad
  m_{\ell} = \min(m(\Lambda_{\ell})) = \min(\{ m_{\ell,k} : k \geq 0\}).
$$  
Also, let
\begin{equation}\label{eq:fell}
  f_{\ell} = \min\{k : m_{\ell,k} < \infty\} \leq \infty,
\end{equation}
where we let $f_{\ell} = \infty$ if $m_{\ell}=\infty$. 

Kolyvagin
proves \cite[Thm.~C]{kolyvagin:structureofsha} that $m_{\ell,0}\geq
m_{\ell,1} \geq m_{\ell,2} \geq \dots $.  

\begin{conjecture}[Kolyvagin's Conjecture $A_\ell$]\label{conj:A}
$m_{\ell} < \infty$.  Equivalently, there exists $\lambda\in \Lambda_{\ell}$ such
that $[P_{\lambda}] \neq 0$.
\end{conjecture}
See \cite{jetchev-lauter-stein} for the first computational evidence
for Conjecture~\ref{conj:A}.  For example, for a specific rank $2$
elliptic curve, that paper shows that $m_3 = m_{3,1} = 0$ and $f_3 =
1$, assuming that the numerical computation of a certain Heegner point
$y_{\lambda}$ was done to sufficient precision. (If the computation
were not done to sufficient precision it is highly likely that we
would haved detected this.)

Conjecture~\ref{conj:A} is quite powerful, as the following theorem
shows.  For an abelian group $A$ of odd order with an action of
complex conjugation, let $A^+$ denote the $+1$ eigenspace for
conjugation and $A^-$ the minus eigenspace, so $A = A^+ \oplus A^-$.
As always, we continue to assume our minimality hypothesis that
$$
r_{\an}(E/\QQ) > r_{\an}(E^D/\QQ) \leq 1.
$$
\begin{theorem}[Kolyvagin]\label{thm:kolythm1}
Let $\ell\in B(E)$, suppose Conjecture~\ref{conj:A} is true for $\ell$,
and let $f = f_{\ell}$. 
For every $k$, let $b_k = \ell^{m_{\ell,k}-m_{\ell,k+1}}$.
Then for every 
$n\geq m_{\ell,f}$, we have
$$
  \Sel^{(\ell^n)}(E/\QQ) = \Sel^{(\ell^n)}(E/K)^+ \ncisom (\ZZ/\ell^n\ZZ)^{f+1} 
     \oplus (\ZZ/b_{f+1}\ZZ)^{2}
     \oplus (\ZZ/b_{f+3}\ZZ)^{2}
     \oplus \cdots
$$
and
$$
  \Sel^{(\ell^n)}(E^D/\QQ) =   \Sel^{(\ell^n)}(E/K)^- \ncisom 
(\ZZ/\ell^n\ZZ)^{h} 
     \oplus (\ZZ/b_f \ZZ)^{2}
     \oplus (\ZZ/b_{f+2} \ZZ)^{2}
     \oplus \cdots
$$
where $h=\rank(E^D(\QQ)) \leq 1$.
\end{theorem}
\begin{proof}
  The leftmost equality in the above two equations is true because
  $\ell$ is odd, and Theorem~1 of \cite{kolyvagin:structure_of_selmer}
  implies both of the rightmost equalities, but possibly with
  $\Sel^{(\ell^n)}(E/K)^+$ and $\Sel^{(\ell^n)}(E/K)^-$ swapped and a
  different value for $h$.  Theorem~1 of
  \cite{kolyvagin:structure_of_selmer} is proved by inductively
  constructing cohomology classes with good properties with respect to
  certain localization homomorphisms.  To finish the proof, we
  establish that these two Selmer groups are not swapped and that
  $h=\rank(E^D(\QQ))$.

  First note that by \cite{kolyvagin:subclass,
    bump-friedberg-hoffstein:nonvanishing}, our hypothesis that
  $r_{\an}(E^D/\QQ) \leq 1$ implies that $r_{\an}(E^D/\QQ) = \rank(E^D(\QQ))$ and
  $\Sha(E^D/\QQ)$ is finite. 

  If $f=0$, then the Heegner point $y_K$ has infinite order, so by
  \cite{gross-zagier} we have $r_{\an}(E/K)=1$ and by
  \cite{kolyvagin:subclass}, $E(K)$ has rank $1$ and $\Sha(E/K)$ is
  finite.  By our minimality hypothesis, we have $r_{\an}(E/\QQ)>
  r_{\an}(E^D/\QQ)$, so $r_{\an}(E/\QQ) = \rank(E(\QQ)) = 1$ and
  $r_{\an}(E^D/\QQ)=\rank(E^D(\QQ)) = 0$.  Thus the two displayed
  Selmer groups $\Sel^{(\ell^n)}(E/K)^{\pm}$ are in the claimed order.
 Moreover, $h=0=\rank(E^D(\QQ))$.

 Next assume $f>0$.  Then one of the two Selmer groups contained
 $(\ZZ/\ell^n\ZZ)^{f+1}$ for arbitrarily large $n$.  Since we know
 that $\Sha(E^D/\QQ)$ is finite and $\rank(E^D(\QQ))\leq 1$ but
 $f+1\geq 2$, the Selmer group that contains $(\ZZ/\ell^n\ZZ)^{f+1}$
 must be $\Sel^{(\ell^n)}(E/K)^+$.  Thus again we see that the two
 displayed Selmer groups are in the claimed order.  Also, again 
$h =  \rank(E^D(\QQ))$ follows.

\qed\end{proof}

\begin{remark}
  Suppose the hypotheses of Theorem~\ref{thm:kolythm1} are satisfied.
  Then comparing the conclusion about the choice of signs in
  Theorem~\ref{thm:kolythm1} with the statement of Theorem~1 in
  \cite{kolyvagin:structure_of_selmer} shows that $f+1 \equiv
  r_{\an}(E/\QQ)\pmod{2}$, which implies the parity conjecture for
  the Selmer group of $E$ at $\ell$.
\end{remark}

\begin{proposition}
  Let $\ell \in B(E)$.  Then $f_{\ell} = \rank(E(\QQ)) - 1$ if and
  only if $\Sha(E/\QQ)(\ell)$ is finite and Conjecture~\ref{conj:A}
  holds for $\ell$.
\end{proposition}
\begin{proof}
  First suppose $f_{\ell} = \rank(E(\QQ)) - 1$.  Then
  $f_{\ell}\neq\infty$, so Conjecture~\ref{conj:A} holds.  To prove
  that $\Sha(E/\QQ)(\ell)$ is finite, use Theorem~\ref{thm:kolythm1}
  and that by our rank hypothesis the image of $E(\QQ)$ in
  $\Sel^{(\ell^n)}(E/\QQ)$ is $(\ZZ/\ell^n\ZZ)^{f+1}$. Thus
  $\Sha(E/\QQ)[\ell^n]$ is a quotient of the $\ell^n$-torsion subgroup
  of the finite group $(\ZZ/b_{f+1}\ZZ)^{2} \oplus
  (\ZZ/b_{f+3}\ZZ)^{2} \oplus \cdots$, so $\Sha(E/\QQ)(\ell)$ is finite.

  Conversely, suppose the $\ell$-primary group $\Sha(E/\QQ)(\ell)$ is
  finite and that Conjecture~\ref{conj:A} holds.    Let $b$ be a positive integer such that $\ell^b
  \Sha(E/\QQ)(\ell) = 0$.  Then the map $\Sel^{(\ell^b)}(E/\QQ) \to
  \Sha(E/\QQ)(\ell)$ is surjective, and for every integer $n\geq b$,
  the map $\Sel^{(\ell^n)}(E/\QQ)[\ell^b] \to \Sha(E/\QQ)(\ell)$ is
  also surjective, since $\Sel^{(\ell^b)}(E/\QQ) \to
  \Sel^{(\ell^n)}(E/\QQ)[\ell^b]$.  
  The image of $\ell^b \Sel^{(\ell^n)}(E/\QQ)$ in $\Sha(E/\QQ)(\ell)$ is trivial.
Since $\ell\in B(E)$ we have
  $E(\QQ)_{\tor}[\ell]=0$, so exactness of the sequence
  $$
   0 \to E(\QQ)/\ell^n E(\QQ) \to \Sel^{(\ell^n)}(E/\QQ) \to \Sha(E/\QQ)(\ell)\to 0
  $$
  implies that $\ell^b \Sel^{(\ell^n)}(E/\QQ) \ncisom
  \ell^b(\ZZ/\ell^{n}\ZZ)^r$, where $r =
  \rank(E(\QQ))$.  On the other hand, if we also choose
  $\ell^b\geq b_{f+1}$, then Theorem~\ref{thm:kolythm1} implies that $
  \ell^b \Sel^{(\ell^n)}(E/\QQ) \ncisom \ell^b (\ZZ/\ell^{n}\ZZ)^{f+1}.$ 
  We conclude that $r=f+1$.  

\qed \end{proof}

 
Kolyvagin's other conjectures involve
$\H^1(K,E[\ell^\infty]) = \displaystyle \varinjlim_{m} \H^1(K,E[\ell^m])$. 

\begin{lemma}\label{lem:conjd}
Suppose $E(K)[\ell]=0$.  Then for every $m\geq 1$, the natural map 
$\H^1(K,E[\ell^m]) \to \H^1(K,E[\ell^\infty])$ is injective.
\end{lemma}
\begin{proof}
  This lemma is of course very well known, but we give a proof for
  completeness.  It suffices to show that for any pair $a,b$ of
  nonnegative integers that the map
\begin{equation}\label{resmn}
\H^1(K,E[\ell^a])\to \H^1(K,E[\ell^{a+b}])
\end{equation}
is injective. 
Taking Galois cohomology of
$0 \to E[\ell^a] \to E[\ell^{a+b}] \to E[\ell^{a+b}]/E[\ell^a] \to 0$
we see that $\H^0(K,E[\ell^{a+b}]/E[\ell^a])$ surjects
onto the kernel of \eqref{resmn}.
We have an exact sequence of Galois modules
$$
 0 \to E[\ell^a] \to E[\ell^{a+b}] \xrightarrow{\ell^a} E[\ell^b] \to 0,
$$
so $\H^0(K,E[\ell^{a+b}]/E[\ell^a]) \isom \H^0(K,E[\ell^b]) = E(K)[\ell^b] = 0$,
since $E(K)[\ell]=0$. 

\qed\end{proof}

We now define Galois cohomology classes associated to the Kolyvagin
points $P_{\lambda}$.  For $\lambda \in \Lambda_{\ell^n}$ with
$\ell\in B(E)$, let $\tau_{\lambda,\ell^n} \in \H^1(K, E[\ell^n])$ be
the image of $P_{\lambda}$ under the map
$$(E(K_{\lambda})/\ell^n E(K_{\lambda}))^{\Gal(K_{\lambda}/K)}
\hookrightarrow 
\H^1(K_{\lambda}, E[\ell^n])^{\Gal(K_{\lambda}/K)}
\isom 
\H^1(K, E[\ell^n]),
$$ 
where the last map is an isomorphism because $\ell \in B(E)$ 
(see, e.g., \cite[\S4]{gross:kolyvagin}).  
Kolyvagin also remarks that one can define Galois cohomology
classes $\tau_{\lambda,\ell^n}$ for $\ell\not\in B(E)$ and all
$\lambda \in \Lambda_{\ell^{k_0+n}}$, where 
$k_0$ is the smallest nonnegative even integer such that 
$\ell^{k_0/2} E(\mathbf{K})(\ell) = 0$ and $\mathbf{K}$ is the compositum of all
$K_{\lambda}$ for $\lambda \in \Lambda$.  Of course, for all $\ell\in B(E)$ we
have $k_0=0$.

%In \cite[pg.~258]{kolyvagin:structure_of_selmer}, Kolyvagin defines
Let $\tau_{\lambda,\ell^n}'$ be the image in $H^1(K,E[\ell^\infty])$ 
of $\tau_{\lambda,\ell^n}$ (note that for the moment we are not
assuming that $\ell\in B(E)$, so the natural map
$\H^1(K,E[\ell^m]) \to \H^1(K,E[\ell^\infty])$ need not be injective).
For any integers $a\geq 0$, $k\geq k_0$ and $n\geq 1$, let
$$
  V_{k,\ell^n}^a = \langle \tau_{\lambda, \ell^n}' : \lambda \in \Lambda_{\ell^{n + k}}^a\rangle
    \subset H^1(K,E[\ell^\infty])
$$
Since $\Lambda_{\ell^{n+k+1}}^a \subset \Lambda_{\ell^{n+k}}^a$, we have
$$V_{0,\ell^n}^a \supset V_{1,\ell^n}^a \supset V_{2,\ell^n}^a \supset \cdots.$$
We have $\ell \tau_{\lambda,\ell^{n+1}} = \tau_{\lambda,\ell^n}$, because
the following diagram commutes, with $G=\Gal(K_{\lambda}/K)$:
%\begin{center}
%\begin{tikzpicture}
% \matrix (m) [matrix of math nodes, row sep=3em,
%    column sep=3em]{
%(E(K_{\lambda})/\ell^{k+n+1}E(K_{\lambda}))^G & \H^1(K_{\lambda},E[\ell^{k+n+1}])^G\\
%(E(K_{\lambda})/\ell^{k+n}E(K_{\lambda}))^G & \H^1(K_{\lambda},E[\ell^{k+n}])^G \\
%};
%\path[right hook->]
%  (m-1-1) edge node[auto] {} (m-1-2)
%  (m-2-1) edge node[auto] {} (m-2-2);
%\path[->]
%  (m-2-1) edge node[auto] {$[\ell]$} (m-1-1)
%  (m-2-2) edge node[auto] {} (m-1-2);
%\end{tikzpicture}
%\end{center}

$$
\xymatrix{
{(E(K_{\lambda})/\ell^{k+n+1}E(K_{\lambda}))^G\,\,}\ar@{^(->}[r] & {\H^1(K_{\lambda},E[\ell^{k+n+1}])^G}\\
{(E(K_{\lambda})/\ell^{k+n}E(K_{\lambda}))^G\,\,}\ar@{^(->}[r]\ar[u]^{[\ell]} & \H^1(K_{\lambda},E[\ell^{k+n}])^G\ar[u] \\
}
$$


Thus $\ell V_{k,\ell^{n+1}}^a \subset V_{k,\ell^n}^a$.

We say $\{\tau_{\lambda, \ell^n}\}$ is a {\em strong nonzero system} if
there exists $a\geq 0$ such that for all $k\geq k_0$ there
exists $n$ such that $V_{k,\ell^n}^a\neq 0$.  In other words,
if one continues the grid of subgroups of $\H^1(K,E[\ell^{\infty}])$
below infinitely far to the right and up in the obvious way, then it
is {\em not} the case that  sufficiently far to
the right every single group is 0.

%\begin{center}
%\begin{tikzpicture}
% \matrix (m) [matrix of math nodes, row sep=0.1em,
%    column sep=2em]{
%   \vdots & \vdots & \vdots  &  \\
% V_{0,\ell^3}^a &  V_{1,\ell^3}^a &  V_{2,\ell^3}^a & \cdots\\
% V_{0,\ell^2}^a &  V_{1,\ell^2}^a &  V_{2,\ell^2}^a & \cdots\\
% V_{0,\ell}^a &  V_{1,\ell}^a &  V_{2,\ell}^a & \cdots\\
%};
%  \path[left hook->]
%  (m-2-4) edge node[auto] {} (m-2-3)
%  (m-2-3) edge node[auto] {} (m-2-2)
%  (m-2-2) edge node[auto] {} (m-2-1)
%  (m-3-4) edge node[auto] {} (m-3-3)
%  (m-3-3) edge node[auto] {} (m-3-2)
%  (m-3-2) edge node[auto] {} (m-3-1)
%  (m-4-4) edge node[auto] {} (m-4-3)
%  (m-4-3) edge node[auto] {} (m-4-2)
%  (m-4-2) edge node[auto] {} (m-4-1)
%;
%\end{tikzpicture}
%\end{center}

$$
\xymatrix@=1.1pc{
   \vdots & \vdots & \vdots  &  \\
 V_{0,\ell^3}^a &  V_{1,\ell^3}^a\ar[l] &  V_{2,\ell^3}^a\ar[l] & \cdots\ar[l]\\
 V_{0,\ell^2}^a &  V_{1,\ell^2}^a\ar[l] &  V_{2,\ell^2}^a\ar[l] & \cdots\ar[l]\\
 V_{0,\ell}^a &  V_{1,\ell}^a\ar[l] &  V_{2,\ell}^a\ar[l] & \cdots\ar[l]\\
}
$$

\begin{conjecture}[Kolyvagin's Conjecture $B_{\ell}$]\label{conj:B}
  $\{\tau_{\lambda, \ell^n}\}$ is a strong nonzero system.
\end{conjecture}

\begin{remark}
  Kolyvagin remarks 
  \cite[pg.~258]{kolyvagin:structure_of_selmer} that if $\ell \in
  B(E)$, then $\{\tau_{\lambda, \ell^n}\}$ is a strong nonzero system
  if and only if there exists $n$ such that $V_{0,\ell^n}^a\neq 0$.
  By Lemma~\ref{lem:conjd}, this is the case if and only if
 some $\tau$ is nonzero.  So
for $\ell\in B(E)$,
Conjectures~\ref{conj:A} is true if and only if Conjecture~\ref{conj:B}
is true.
\end{remark}

The following conjecture is motivated by Theorem~\ref{thm:kolythm1}
and the conjecture that $\Sha(E/K)$ is finite. 
\begin{conjecture}[Kolyvagin's Conjecture C]\label{conj:C} 
  The set of primes $\ell$ such that $m_{\ell}\neq 0$ is finite.
\end{conjecture}

Let $r_{\an} = \ord_{s=1} L(E,s)$, and let $\varepsilon = (-1)^{r_{\an}-1}$. 
For any module $A$ with an action of complex conjugation $\sigma$,
and $\nu\in\{0,1\}$, let $A^\nu = (1-(-1)^\nu \varepsilon \sigma)A$.

\begin{conjecture}[Kolyvagin's Conjecture $D_{\ell}$]\label{conj:D} 
  There exists $\nu\in\{0,1\}$ and a subgroup $V\subset
  (E(K)/E(K)_{\tor})^{\nu}$ such that $1\leq \rank(V) \equiv \nu
  \pmod{2}$ and for all $n\geq 1$ and all sufficiently large $k$, one has
  $$
 V_{k,\ell^n}^a \equiv V\pmod{\ell^n (E(K)_{/\tor})},
$$ 
   where $a = \rank(V)-1$.
\end{conjecture}

The following conjecture is the natural generalization to higher rank
of the hypothesis when $r_{\an}(E/\QQ)=1$ that the Hegner point $y_K$
has infinite order.
\begin{conjecture}[Kolyvagin's Conjecture D]\label{conj:univD}
  There exists a single subgroup $V$ of $E(K)$ such that
  Conjecture~\ref{conj:D} holds simultaneously for all $\ell$ with
  that $V$.
\end{conjecture}

Conjecture~\ref{conj:D} has numerous consequences.  Much of the
following proposition is implicitly stated without any proofs in
\cite[pg.~258--259]{kolyvagin:structure_of_selmer}, so we give
complete proofs below.
\begin{proposition}\label{prop:conjD}
Assume our running minimality hypothesis that 
$r_{\an}(E/\QQ) > r_{\an}(E^D/\QQ) \leq 1$. 
Suppose Conjecture~\ref{conj:D} is true for $\ell\in B(E)$ and
let $f=f_{\ell}$. 
Then 
\begin{enumerate}
\item\label{prop:D:1}
 $
(E(K)/E(K)_{\tor})^{\nu} = (E(K)/E(K)_{\tor})^+,
$ 
\item\label{prop:D:2}
$a = f$, 
\item\label{prop:D:3}
 $\rank(E(\QQ)) = f + 1$, 
\item\label{prop:D:4} 
$\Sha(E/K)(\ell)$ is finite,
\item\label{prop:D:5} 
$r_{\an}(E/\QQ) \equiv \rank(E(\QQ))\pmod{2}$, and
\item\label{prop:D:6} 
 $V\tensor\ZZ_{\ell} = \ell^{m_{f}} E(\QQ)\tensor\ZZ_{\ell}$.
\end{enumerate}
\end{proposition}
\begin{proof}
By Conjecture~\ref{conj:D}, there exists $\nu\in\{0,1\}$ and a subgroup
$V\subset (E(K)/E(K)_{\tor})^{\nu}$ such that $1\leq \rank(V) \equiv \nu\pmod{2}$
and for all $n>0$ and all sufficiently large $k$ we have
$$
  V_{k,\ell^n}^a \equiv V \pmod{ \ell^n E(K)_{\tor} },
$$
where $a = \rank(V) - 1$.  

If $\rank(V)=1$, then $a=0$, so $V_{k,\ell^n}^0 \neq 0$ for some $k$,
so since $f$ is the smallest integer such that $V_{k,\ell^n}^f \neq 0$,
this implies that $f=0$ giving Part~\ref{prop:D:2}; thus the Heegner
point $y_K$ has infinite order and $r_{\an}(E/K) = 1$.  Since
$r_{\an}(E/\QQ) > r_{\an}(E^D/\QQ)$, we have $r_{\an}(E/\QQ)=1$ and
$r_{\an}(E^D/\QQ)=0$, so
Parts~\ref{prop:D:1},\ref{prop:D:3},\ref{prop:D:4}, \ref{prop:D:5}
follows.  Finally, Part~\ref{prop:D:6} follows since $V_{k,\ell^n}^0$
is just the image of the Heegner point $y_K$ under the connecting
homomorphism, and $\ord_{\ell}(y_K) = m_f$.

Next assume that $\rank(V)>1$.  By our minimality hypothesis,
$r_{\an}(E^D/\QQ)\leq 1$, so $\rank(E^D(\QQ))\leq 1$, hence
$V\not\subset (E(K)/E(K)_{\tor})^-$, so $V\subset
(E(K)/E(K)_{\tor})^+$, which proves Part~\ref{prop:D:1}.  We have
$f\leq a$ since $V_{k,\ell^n}^a\neq 0$ for some $k\geq 0$.  Also,
since $f<\infty$, Theorem~\ref{thm:kolythm1} implies that
$\rank(E(\QQ)) \leq f+1$.  Since $\rank((E(K)/E(K)_{\tor})^+) =
\rank(E(\QQ))$, we have $$a+1 = \rank(V) \leq \rank(E(\QQ)) \leq f+1
\leq a+1.$$ We conclude that the above inequalities are equalities, so
$a=f$ which proves Part~\ref{prop:D:2}, and $\rank(E(\QQ)) = f+1$,
which proves Part~\ref{prop:D:3}.  Also because $\rank(E(\QQ))=f+1$,
Theorem~\ref{thm:kolythm1} implies that $\Sha(E/\QQ)(\ell)$ is finite,
so since $\Sha(E^D/\QQ))$ is also finite, Part~\ref{prop:D:4} is
true.  Considering the definition of the $A^{\nu}$ before the
statement of Conjecture~\ref{conj:D}, we see that
$1-(-1)^{\nu}(-1)^{r_{\an}-1}\sigma = 1 + \sigma$, so $\nu \equiv
r_{\an} \pmod{2}$.  Since part of Conjecture~\ref{conj:D} is that
$\rank(V)\equiv \nu\pmod{2}$, and we proved that
$\rank(V)=\rank(E(\QQ))$, we conclude that $r_{\an}\equiv
\rank(E(\QQ))\pmod{2}$, which is Part~\ref{prop:D:5}.  By
\cite[Thm.~3]{kolyvagin:structure_of_selmer},
%\footnote{It is not
%  completely clear to the author that Theorem~3 of
%  \cite{kolyvagin:structure_of_selmer} implies that $\delta(\ell^{m_f}
%  E(\QQ)) \subset V_{k,\ell^n}^f$ for {\em all} $k\geq m_f$.  Kolyvagin
%  explicitly states this in the middle of
%  \cite[pg.~258]{kolyvagin:structure_of_selmer}.  The point of
%  Proposition~\ref{prop:conjD} is to lay out the {\em consequences} of Kolyvagin's
%  conjecture that we will use to conditionally deduce properties of
%  certain groups $W_p$ later, so if the reader also doesn't consider
%  the statement sufficiently justified, then the reader can simply
%  consider this part of the proposition part of Kolyvagin's
%  conjecture.} 
for all $k\geq m_f$ the subgroup
$V_{k,\ell^n}^f\subset \H^1(K,E[\ell^\infty])$ contains
$(\ell^{m_f}\ZZ/\ell^n\ZZ)^{f+1} = \delta(\ell^{m_f} E(\QQ))$, so
$\ell^{m_f} E(\QQ) \tensor \ZZ_{\ell} \subset V \tensor \ZZ_{\ell}$.
On the other hand, by definition of $m_f$, every cohomology class
$\tau_{\lambda, \ell^n}$ is contained in $\ell^{m_f}
\H^1(K,E[\ell^n])$.  Thus $\delta(V) \subset \ell^{m_f}
\H^1(K,E[\ell^\infty])$, so $V\subset \ell^{m_f} E(\QQ)$.  This proves
Part~\ref{prop:D:6}.

\qed\end{proof}



Recall from Section~\ref{sec:ggz} that $c$ is the Manin constant of
$E$ and the $c_q$ are the Tamagawa numbers of $E$.  We make the
following new refinement of Kolyvagin's Conjecture~\ref{conj:C}. 
\begin{conjecture}\label{conj:refined}
We have $m_{\ell} = \ord_{\ell}(c \cdot \prod_{q\mid N} c_q)$.
\end{conjecture}
Theorem~\ref{thm:refined_evidence} and Theorem~\ref{thm} below serve
as our motivation to make Conjecture~\ref{conj:refined}.  In
particular, Kolyvagin proved that at primes $\ell\in B(E)$,
Conjecture~\ref{conj:refined} is equivalent to \cite[Conj~2.2, pg
311]{gross-zagier} in the special case when $E$ has analytic rank $1$
over $K$.

\section{Mod~$p$ Kolyvagin Points and Kolyvagin Subgroups} 
\label{sec:kolysubgroups} 

As always, we assume $E$ is an elliptic curve over $\QQ$, that $K$ is
a quadratic imaginary field satisfying the Heegner hypothesis, and $p$
is a prime that is inert in $K$.  The Heegner hypothesis implies that
the primes of bad reduction for $E$ split in $K$, so $p$ must be a
prime of good reduction.  For each such prime, we define a
finite-index subgroup $W_p$ of $E(K)$.  We do this by extending
Kolyvagin's construction of points $P_{\lambda}$ to obtain a new
well-defined construction of elements of the quotient group
$$
E(\FF_p)/(p+1) = E(\FF_p)/(p+1) E(\FF_p)
$$ 
for any inert prime $p$.  Thus this section takes Kolyvagin's
definition of points $P_{\lambda}$ one step further to define elements
of $E(\FF_p)/(p+1)$.  We first compute the structure of the odd part
of the group $E(\FF_p)/(p+1)$ for any good prime $p$. We then use
properties of splitting of primes in certain ring class fields to
define the canonical reduction $R_{p,\lambda}\in E(\FF_p)/(p+1)$ of
the Kolyvagin points $P_{\lambda}$, and consider the subgroup $X_p$ of
$E(\FF_p)/(p+1)$ generated by the $R_{p,\lambda}$ for certain
$\lambda$.  We then define $W_p$ to be the inverse image of~$X_p$
and finish with some results about the structure of $W_p$.

If $A$ is a finite abelian group, the {\em odd part} of $A$ is the
subgroup of $A$ of all elements of odd order, and if $n$ is an
integer, the odd part of $n$ is $n/2^{\ord_2(n)}$.
\begin{lemma}\label{lem:cyclic}
The odd part of $E(\FF_p)/(p+1)$ is cyclic of order the odd part of $\gcd(p+1,a_p)$.
\end{lemma}
\begin{proof}
  Suppose $\ell$ is an odd prime divisor of $\#(E(\FF_p)/(p+1))$.  If
  the $\ell$-primary subgroup of $E(\FF_p)/(p+1)$ is not cyclic, then
  since $\ell\neq p$ we have $E(\FF_p)[\ell] \ncisom (\ZZ/\ell\ZZ)^2$.
  The Weil pairing induces an isomorphism of Galois modules
  $\bigwedge^2 E[\ell] \isom \mu_{\ell}$ and $E[\ell] \subset
  E(\FF_p)$, so $\mu_{\ell}\subset \FF_p^*$, hence $\ell\mid (p-1)$.
  Since $\ell$ divides $\#(E(\FF_p)/(p+1))$ and $\ell$ is prime, we
  have $\ell\mid (p+1)$, so $\ell\mid\gcd(p-1,p+1) = 2$, a
  contradiction, since $\ell$ is odd.

  The group $E(\FF_p)$ has order $p+1-a_p$, and we just proved above
  that $E(\FF_p)(\ell)$ is cyclic for any odd prime divisor $\ell$ of $p+1$.  Thus the
  quotient $\ell$-primary group $(E(\FF_p)/(p+1))(\ell) =
  (E(\FF_p)(\ell))/(p+1)$ has order $\ell^{m}$, where
   $$m = \ord_{\ell}(\gcd(p+1,\#E(\FF_p)))
         =  \ord_{\ell}(\gcd(p+1,p+1-a_p))
         =  \ord_{\ell}(\gcd(p+1,a_p)).
  $$
  Taking the product over all odd primes $\ell$, shows that the odd
  part of $E(\FF_p)/(p+1)$ has order the odd part of $\gcd(p+1,a_p)$.
\qed\end{proof}
\begin{remark}
\begin{enumerate}
\item Lemma~\ref{lem:cyclic} is true even if $p$ is a good prime that
  is not inert in $K$ (in fact, the lemma and proof have nothing to do
  with $K$).
\item Lemma~\ref{lem:cyclic} is false if we do not restrict to odd parts.  For example,
if $E$ is $y^2=x^3 - x$ and $p=3$, then $E(\FF_3) \ncisom (\ZZ/2\ZZ)^2$, so
$E(\FF_3)/4 \ncisom (\ZZ/2\ZZ)^2$ is not cyclic. 
\item For every prime $\ell$, there exists infinitely many primes $p$ such that
$E(\FF_p)(\ell)$ is not cyclic.  Indeed, by the Chebotarev density theorem there
are infinitely many $p$ that split completely in the field $\QQ(E[\ell])$,
and for these $p$ we have $(\ZZ/\ell\ZZ)^2\subset E(\FF_p)$.
\end{enumerate}
\end{remark}

\begin{lemma}\label{lem:splits}
  If $p$ is inert in $K$ and does not divide $\lambda$, then the prime
  ideal $p\OO_K$ of $K$ splits completely in $K_{\lambda}$.  In
  particular, if $p\in \Lambda^1_{\ell^n}$ and $\lambda \in
  \Lambda_{\ell^n}$ with $p\nmid \lambda$, then $p\OO_K$ splits
  completely in $K_{\lambda}$.
\end{lemma}
\begin{proof}
(Compare line $-3$ on page 103 of \cite{kolyvagin:structureofsha}.)
Since $p$ is inert, the ideal $p\OO_K$ is a prime principal ideal
of $\OO_K$, hence splits completely in the Hilbert class field $K_1$.
As explained in \cite[pg.~238]{gross:kolyvagin}, class field theory
identifies $\Gal(K_{\lambda}/K_1)$ with $C = (\OO_K/\lambda\OO_K)^* /
(\ZZ/\lambda\ZZ)^*$.  The image of $p$ is trivial in $C$, so the
Frobenius element attached to $p\OO_K$ is trivial, hence $p\OO_K$
splits completely in the ring of integers of $K_{\lambda}$, as
claimed.  \qed \end{proof}

Define the reduction map $E(K)\to E(\FF_{p^2})$ by reducing the
N\'eron model $\mathcal{E}$ of $E$ over $\OO_K$ modulo $p\OO_K$, and
using the natural maps $E(K)\isom\mathcal{E}(\OO_K) \to
\mathcal{E}_{\FF_{p^2}}(\FF_{p^2}) \isom E(\FF_{p^2})$.  Let $\pi_p:
E(K) \to E(\FF_p)/(p+1)$ be the composition of reduction modulo the
prime ideal $p\OO_K$ with $\Tr_{\FF_{p^2}/\FF_p}:E(\FF_{p^2})\to
E(\FF_p)$ followed by quotienting out by the subgroup $(p+1)E(\FF_p)$.
Fix a {\em choice} $\wp$ of prime ideal of $K_{\lambda}$ over
$p\OO_K$.  Extend $\pi_p$ to a map $\pi_{\wp} : E(K_{\lambda}) \to
E(\FF_{p})/(p+1)$ by quotienting out by $\wp$, as illustrated in the
following diagram:
%\begin{center}
%\begin{tikzpicture}
% \matrix (m) [matrix of math nodes, row sep=2em,
%    column sep=2em]{
%E(K_{\lambda}) &          & \\
%E(K) & \mathcal{E}(\OO_K) & E(\FF_{p^2}) \\
%     &                    & E(\FF_p) \\ 
%     &                    & E(\FF_p)/(p+1) \\};
%\path[right hook->]
%   (m-2-1) edge node[auto] {} (m-1-1);
%\path[->]
%   (m-1-1) edge node[auto] {$\mod{\wp}$} (m-2-3)
%   (m-2-1) edge node[auto] {$\isom$} (m-2-2)
%   (m-2-2) edge node[auto] {} (m-2-3)
%   (m-2-3) edge node[auto] {trace} (m-3-3)
%   (m-3-3) edge node[auto] {} (m-4-3)
%   (m-2-1) edge [densely dotted] node[below] {$\pi_p$} (m-4-3)
%   (m-1-1) edge [densely dotted, bend right=90] node[below] {$\pi_{\wp}$} (m-4-3);
%\end{tikzpicture}
%\end{center}

$$
\xymatrix{
E(K_{\lambda})\ar[drr]^{\mod{\wp}} &          & \\
E(K)\ar[u]\ar[r]^{\isom}\ar@{.>}[ddrr]_{\pi_p} & \mathcal{E}(\OO_K)\ar[r] & E(\FF_{p^2})\ar[d]^{{\rm trace}} \\
     &                    & E(\FF_p)\ar[d] \\ 
     &                    & E(\FF_p)/(p+1)\\
}
$$

For each $\ell\mid (p+1)$, let 
$v_{\ell} = \ord_{\ell}(\gcd(a_p, p+1))$,
and define
$\pi_{\wp,\ell} : E(K_{\lambda}) \to (E(\FF_{p})/(p+1)) (\ell)$
by
$$
  \pi_{\wp,\ell}(S) = \pi_{\wp}\left(\frac{p+1}{\ell^{v_\ell}} S \right).
$$

We now study how the homomorphism $\pi_{\wp,\ell}$ depends on our
choice of prime of $\wp$ over $p\OO_K$. 

\begin{proposition}\label{prop:piplambdaell}
The map $\pi_{\wp,\ell}$ induces a well-defined (independent of choice
of $\wp$) homomorphism
$$
  \vartheta : \left(E(K_{\lambda})/ \ell^{v_{\ell}} E(K_{\lambda})\right)^{\Gal(K_{\lambda}/K)}
    \to E(\FF_p) / (p+1).
$$
\end{proposition}
\begin{proof}
Let $[S] \in \left(E(K_{\lambda})/ \ell^{v_{\ell}} E(K_{\lambda})\right)^{\Gal(K_{\lambda}/K)}$
with $S\in E(K_{\lambda})$.
If $\wp'$ is another prime of $K_{\lambda}$ over $p\OO_K$, then because the Galois group
acts transitively on the primes over a given prime, there
is $\sigma\in \Gal(K_{\lambda}/K)$ such that
$\pi_{\wp',\ell}(S) = \pi_{\wp,\ell}(\sigma(S))$.
Since $[S]$ is $\Gal(K_{\lambda}/K)$-equivariant, we have
$
  \sigma(S) = S + \ell^{v_{\ell}} \cdot Q,
$
for some $Q\in E(K_{\lambda})$, so
\begin{align*}
\vartheta([\sigma(S)]) &= \pi_{\wp,\ell}(\sigma(S))\\
& =  \pi_{\wp}\left(\frac{p+1}{\ell^{v_\ell}} \sigma(S) \right)\\
 &=
 \pi_{\wp}\left(\frac{p+1}{\ell^{v_\ell}} S \right) + \pi_{\wp}((p+1)Q)\\
& = 
\pi_{\wp,\ell}(S) + 0 = \vartheta([S]),
\end{align*}
where $\pi_{\wp}((p+1)Q) = (p+1)\pi_{\wp}(Q)$ is $0$, since the
group $E(\FF_p)/(p+1)$ is killed by $p+1$. 
\qed\end{proof}

By Proposition~\ref{prop:invariant}, $[P_{\lambda}]$ is in the
domain of the homomorphism $\vartheta$ of Proposition~\ref{prop:piplambdaell}.
\begin{definition}[Mod~$p$ Kolyvagin Point]
The {\em mod~$p$ Kolyvagin point} associated to $p\in \Lambda_{\ell^n}^1$ and
$\lambda\in\Lambda_{\ell^n}$ is
$$
  R_{p,\lambda} = \vartheta ( [P_{\lambda}] ) \in E(\FF_p)/(p+1),
$$
where $\vartheta$ is as in Proposition~\ref{prop:piplambdaell}.
\end{definition}
As above, let $v_{\ell} = \ord_{\ell}(\gcd(a_p, p+1))$. For each $k\geq 0$, let
\begin{equation}\label{eqn:Xp}
X_{k,p} = \left\langle R_{p,\lambda} : 
    \lambda \in \bigcup_{\ell} \Lambda_{\ell^{v_{\ell}+k}}^{f_{\ell}} \right\rangle
\subset E(\FF_p)/(p+1)
\end{equation}
be the subgroup generated by all mod~$p$ Kolyvagin points associated
to $\lambda$ that are a product of $f_{\ell}$ primes, where $f_{\ell}$
is from Equation~\eqref{eq:fell}.  Note that the subscript of
$\Lambda$ in \eqref{eqn:Xp} is $\ell^{v_\ell+k}$, and we take the
union over {\em all $\ell$} thus obtaining a subgroup $X_{k,p}$ that
need not be $\ell$-primary for any~$\ell$, despite $R_{p,\lambda}$
being $\ell$-primary.  Let
$$
  X_p = \bigcap_{k\geq 0} X_{k,p}.
$$
Let $W_{k,p}$ be the inverse image of $X_{k,p}$ under the
map $\pi_p$:
$$
  W_{k,p} = \pi_p^{-1}(X_{k,p}) \subset E(K),
$$
and 
$$
  W_{p} = \pi_p^{-1}(X_{p}) \subset E(K).
$$
Since $E(\FF_p)/(p+1)$ is finite, $W_{k,p}$ and $W_p$ have finite index in $E(K)$;
also, by Lemma~\ref{lem:cyclic}, the odd part of this index divides
$\gcd(p+1,a_p)$ .  

\begin{remark}\label{rem:kertr}
  Note that $E^D(\QQ)$ is in the kernel of the trace map, hence in the
  kernel of $\pi_p$, so $E^D(\QQ) \subset W_p$.  Thus it is possible
  that $W_p$ contains torsion, hence $W_p$ in general need not be a
  Gross-Zagier subgroup as in Definition~\ref{defn:ggz}.  In a future
  paper, we intend to give a more refined definition of a sequence of
  groups $W_p^a$, for each $a\geq 0$, which better accounts for
  torsion.  We would then search for a Gross-Zagier style formula for
  each group $W_p^a$ for $a \leq f+1$, in order to more closely relate
  $r_{\an}(E/\QQ)$ to $f+1$.
\end{remark}


%\begin{remark}
%  We could define $X_p^a$ for any $a$ (not just $a=f$) as above,
%  except if $a\not\equiv f\pmod{2}$, we would project to $E'(\FF_p)$,
%  where $E'$ is the quadratic twist of $E_{\FF_p}$.
%\end{remark}


\section{Controlling the Reduction Map}\label{sec:reduction}
The main result of this section is a proof that under certain
hypothesis, if a point $Q$ has infinite order and $n$ is a positive
integer, then there are infinitely many primes $p$ such that the image
of $Q$ in $E(\FF_{p^2})/(p+1)$ has order divisible by $n$.  We prove
this using Galois cohomology and by converting a condition on
$\ell$-divisibility of points into a Chebotarev condition.  We will
use this result later to study the maximum index $[E(K):W_p]$ that can
occur and prove a generalized Gross-Zagier formula for such $W_p$.

Let $E$, $K$, etc., be as above, and let $\ell\in B(E)$, where $B(E)$ is
the set of primes defined on page~\pageref{defn:be}.
Suppose $Q\in E(K)$ has infinite order, and let $n$ be an odd positive integer. 
Suppose that for each prime $\ell\mid n$, the set of cardinalities
$\{\#\H^1(K(E[\ell^j])/K,E[\ell^j]) : j \geq 1\}$ is bounded.   This
hypothesis is satisfied if $\ell\in B(E)$, since then
$\H^1(K(E[\ell^j])/K,E[\ell^j]) = 0$ for all $j$
(see \cite[pg.~241]{gross:kolyvagin} and \cite[Prop.~5.2]{bsdalg1}). 
\begin{proposition}~\label{prop:bigord}
Let $Q$ and $n$ be as above. 
Let $S$ be the set of primes $p$ such that $p$
is inert in $K$, $p$ splits completely in $K(E[n])/K$, 
and the image of $Q$ in $E(\FF_{p^2})/(p+1)E(\FF_{p^2})$  has order divisible by~$n$.
Then $S$ has positive (Dirichlet) density.
\end{proposition}
\begin{proof}
 Let $m = \prod \ell_i^{e_i}$ with $\ell_i$
  the distinct primes that divide $n$, and $e_i$ any positive integers,
which we will fix later in the argument.  Fix any $i$, and let $L =
  K(E[\prod_{j\neq i}\ell_j])$, which is a Galois extension of $K$.
  Define homomorphisms $\Psi_i$, $f$, $g$, and $h$ as in the following
  commutative diagram:
%\begin{center}
%\begin{tikzpicture}
% \matrix (m) [matrix of math nodes, row sep=3em,
%    column sep=3em]{
%   E(K(E[m]))/\ell_i^{e_i} E(K(E[m])) &   \H^1(K(E[m]),E[\ell_i^{e_i}]) \\
%        &  \H^1(L(E[\ell_i^{e_i}]),E[\ell_i^{e_i}])\\
%        &  \H^1(K(E[\ell_i^{e_i}]),E[\ell_i^{e_i}])\\
%     E(K)/\ell_i^{e_i} E(K)     &   \H^1(K, E[\ell_i^{e_i}])\\};
%\path[right hook->]
%   (m-1-1) edge node[auto] {} (m-1-2)
%   (m-4-1) edge node[auto] {} (m-4-2);
%\path[->]
%   (m-4-2) edge node[auto] {$g$} (m-3-2)
%   (m-3-2) edge node[auto] {$h$} (m-2-2)
%   (m-2-2) edge node[auto] {$f$} (m-1-2)
%   (m-4-1) edge node[auto] {} (m-1-1)
%   (m-4-1) edge [bend left=15, densely dotted] node[auto] {$\Psi_i$} (m-1-2);
%\end{tikzpicture}
%\end{center}

$$
\xymatrix{
   E(K(E[m]))/\ell_i^{e_i} E(K(E[m]))\,\,\ar@{^(->}[r] &   \H^1(K(E[m]),E[\ell_i^{e_i}]) \\
        &  \H^1(L(E[\ell_i^{e_i}]),E[\ell_i^{e_i}])\ar[u]^f\\
        &  \H^1(K(E[\ell_i^{e_i}]),E[\ell_i^{e_i}])\ar[u]^h\\
     E(K)/\ell_i^{e_i} E(K)\ar[uuu]\ar@{.>}[uuur]^{\psi_i}\,\ar@{^(->}[r]    &   \H^1(K, E[\ell_i^{e_i}])\ar[u]^g
}
$$


The horizontal maps above are induced by the short exact sequence
coming from multiplication by $\ell_i^{e_i}$, and the vertical maps on
the right are the restriction maps.  The diagram commutes so
the order of the image of $Q$ in $E(K(E[m]))/\ell_i^{e_i} E(K(E[m]))$
is the same as the order of $\Psi_i(Q)$.

By hypothesis and the inflation restriction sequence\, the cardinality
of $\ker(g)$ is bounded independently of $i$ and $e_i$.  Also, $[L:K]$
depends only on the set of prime divisors $\ell_i$ of $n$, not their
exponents, so
$$
 \#\ker(h) = \#\H^1(L(E[\ell_i^{e_i}])/K(E[\ell_i^{e_i}]), E[\ell_i^{e_i}])
    = \#\Hom(\Gal(L(E[\ell_i^{e_i}])/K(E[\ell_i^{e_i}])), E[\ell_i^{e_i}])
$$
is also bounded independent of $e_i$, because every homomorphism
has image in the fixed subset $E[\ell_i^{d}]$, where 
$d$ is the exponent of the group $\Gal(L/K)$.    
Finally, the map $f$ is injective, since 
$$
 \ker(f) \isom \H^1(K(E[m])/L(E[\ell_i^{e_i}]), E[\ell_i^{e_i}])
$$
and
$\#\Gal(K(E[m])/L(E[\ell_i^{e_i}]))$ is divisible only
by the primes $\ell_j$ for $j\neq i$ and these are
all coprime to $\#E[\ell_i^{e_i}]=\ell_i^{2e_i}$.   We conclude that there is
an integer $b$ such that $\#\ker(\Psi_i)\leq \ell_i^b$, and
this bound holds no matter how we increase the numbers
$e_i$ and $e_j$ (for all $j$). 

The above proof that $\ker(\Psi_i)$ is uniformly bounded is completely
general. See Remark~\ref{rmk:otherarg} for a sketch of an alternative
proof of this bound in the special case when $\ell\in B(E)$ for all
$\ell\mid n$, which is the only case we will use in this paper.

Because $\ker(\Psi_i)$ is uniformly bounded independent of our choice
of $e_i$, for each $i$, we can choose $e_i$ large enough so that
$\Psi_i(Q)$ has order divisible by $\ell_i^{\ord_{\ell_i}(n)}$.  Then
for each $i$, let $d_i$ be maximal such that $\ell_i^{d_i}$ divides
$Q$ in $E(K(E[m]))$.  Note that $d_i < e_i$ for each $i$, since
$\Psi_i(Q)\neq 0$ and $\Psi(Q)$ is an element of a group that is
killed by $\ell_i^{e_i}$.  Since $m=\prod \ell_i^{e_i}$, we have
$\ell_i^{d_i+1}\mid m$, so
$$M_i = K\left(E[m], \,\,\frac{1}{\ell_i^{d_i+1}} Q\right)$$
does not depend on the choice of $\ell_i^{d_i+1}$th root of $Q$,
is a Galois extension of $K(E[m])$, and
$[M_i:K(E[m])]$ is a nontrivial power of $\ell_i$.
Thus the $M_i$ for all $i$ are linearly disjoint 
as extensions of $K(E[m])$.

Let $M$ be the compositum of the fields $M_i$ defined above.  Since
the $M_i$ are linearly disjoint nontrivial extensions of $K(E[m])$,
there exists an automorphism $\sigma \in \Gal(M/\QQ)$ such that
$\sigma|_{K(E([m]))}$ is complex conjugation, and $\sigma|_{M_i}$ has
order divisible by $\ell_i$ for each $i$.  By the Chebotarev density
theorem, there is a positive density of primes $p\in\ZZ$ that are
unramified in $M$ and have Frobenius the class of $\sigma$.  Such
primes are inert in $K$ since complex conjugation acts nontrivially
on $K$, split completely in $K(E[m])/K$ since complex conjugation
has order $2$, and each prime over $p$ in $K(E[m])$ does not split
completely in any of the extensions $M_i/K(E[m])$ since
$[\Frob_p]|_{M_i}=\sigma|_{M_i}$ has order divisible by $\ell_i>2$.
Note that this is the only place in the argument where we use that $n$ is
odd.

Let $p$ be any prime as in the previous paragraph.  We have 
$$E(\FF_{p^2})/ \ell_i^{e_i} E(\FF_{p^2}) \isom (\ZZ/\ell_i^{e_i}\ZZ)^2$$ 
since $p\OO_K$ splits completely in $K(E[m])$ and $\ell_i^{e_i}\mid m$.
Also, the Frobenius condition implies that the primes of $M_i$ over
$p\OO_K$ do not have residue class degree $1$, so since
$M_i$ is generated by any choice of $\frac{1}{\ell^{d_i+1}}Q$, the
reduction $\overline{Q}$ of $Q$ modulo any prime over $p\OO_K$ is not
divisible by $\ell_i^{d_i+1}$ in $E(\FF_{p^2})$.  Note that
$\ell_i^{d_i}$ divides $\overline{Q}$, because the prime $p\OO_K$ splits
completely in $K(E[m])/K$ and $\ell_i^{d_i}$ divides $Q$ in $K(E[m])$,
so $d_i$ is the largest integer such that $\ell_i^{d_i}$ divides the
image of $\overline{Q}$ in $E(\FF_{p^2})$.   We
conclude that for each $i$ the image of $Q$ in $E(\FF_{p^2})/
\ell_i^{e_i} E(\FF_{p^2})$ has order the same as the order of
$\Psi_i(Q)$.

By hypothesis, $e_i\geq \ord_{\ell_i}(n)$
and $\Psi_i(Q)$ has order divisible by $\ell_i^{\ord_{\ell_i}(n)}$ for each $i$,
so the image of $Q$ in $E(\FF_{p^2})/ m E(\FF_{p^2})$ has
order divisible by $n$. 
For any such $p$, we also have that the characteristic
polynomial of the class of $\Frob_p$ in $\Gal(\QQ(E[m])/\QQ)$ acting on $E[m]$ is
$x^2 - a_p x + p \pmod{m}$.  On the other hand, since $[\Frob_p]$ on $E[m]$
is the class of complex conjugation and complex conjugation acts nontrivially (since
$m$ is odd) hence
has characteristic polynomial $x^2-1$, we have
$x^2 - a_p x + p \equiv x^2 - 1 \pmod{m}$. 
Thus $m \mid (p+1)$, so the image of $Q$ in 
$E(\FF_{p^2})/ (p+1) E(\FF_{p^2})$ also has order divisible by $n$,
which completes the proof.
\qed\end{proof}
% \vspace{2in}
% On the other hand, since $p$ is inert in $K$ and splits completely
% in $E[\ell_i^{e_i}]$, we have
% $$
%   E(\FF_{p^2})[\ell_i^{e_i}] \isom (\ZZ/\ell_i^{e_i}\ZZ)^2,
% $$
% so 
% $$
%   E(\FF_{p^2})/ \ell_i^{e_i}E(\FF_{p^2}) \isom (\ZZ/\ell_i^{e_i}\ZZ)^2.
% $$
%   Since this is true for all $i$,
% we conclude that the image of $Q$ in $E(\FF_{p^2})$ has order 
% divisible by $n$. 


\begin{remark}\label{rmk:bigord}
  Proposition~\ref{prop:bigord} is analogous to the statement that if
  $x,n\in\ZZ$ with $\gcd(n,x)=1$ and $\QQ(\zeta_n, \sqrt[n]{x})$ is an
  extension of $\QQ(\zeta_n)$ of degree $n$, then there exist a
  positive density of primes $p$ such that the multiplicative order of
  $x$ modulo $p$ is divisible by $n$. The proof of this statement
  resembles the proof of Proposition~\ref{prop:bigord}, except we work
  with the field $\QQ(\zeta_n, \sqrt[n]{x})$.  The idea of the proof
  of Proposition~\ref{prop:bigord} is well-known to experts who study
  questions such as the Lang-Trotter conjecture about reduction of
  points on elliptic curves.
\end{remark}

\begin{remark}\label{rmk:otherarg}
  If for every prime $\ell\mid n$ we have $\ell \in B(E)$, we can
  alternatively use that $K(E[\ell_1^{\infty}])$ and
  $K(E[\ell_2^{\infty}])$ are linearly disjoint for distinct odd
  primes $\ell_1$ and $\ell_2$ in $B(E)$ to give a different proof that
  the maps $\Psi_i$ have uniformly bounded kernel in
  Proposition~\ref{prop:bigord}.  In that case we have that
  $\Gal(K(E[n])/K) \ncisom \GL_2(\ZZ/n\ZZ)$, so
$$
  \ker\Bigl(\H^1(K,E[n]) \to \H^1(K(E[n]),E[n])\Bigr) \isom 
   \H^1(K(E[n])/K, E[n]) = \H^1(\GL_2(\ZZ/n\ZZ), (\ZZ/n\ZZ)^2) = 0,
$$
where the last group is $0$ by a standard group cohomology argument
(see, e.g., \cite[\S5.1]{stein:index}).  This implies that the maps
$\Psi_i$ are all injective.  The linear disjointness of
$K(E[\ell_1^{\infty}])$ and $K(E[\ell_2^{\infty}])$ for the distinct
{\em odd} primes $\ell_1$ and $\ell_2$ follows by a Galois theory
argument using the structure of $\GL_2(\ZZ/\ell^n\ZZ)$.  We 
thank R.~Greenberg for this observation.

%% FROM RALPH
% Suppose that p and q are distinct odd primes, that E is an elliptic
% curve/Q, and that the representations giving the Galois action on the
% Tate module for E for p and for q are surjective.  Let K_p =
% Q(E[p^{infty}]) and K_q = Q(E[q^{infty}]). Suppose that K_p and K_q
% have a nontrivial intersection. Then there exists a Galois extension K
% contained in K_p and K_q such that G = Gal(K/Q) is a simple
% group. Thus G is a simple quotient of both GL_2(Z_p) and GL_2(Z_q).
% Assume first that G is nonabelian. Then the images of SL_2(Z_p) and
% SL_2(Z_q) in G will be nontrivial. Those images will be normal
% subgroups of G. Hence those images must be all of G. That is, G is a
% simple quotient of both SL_2(Z_p) and SL_2(Z_q).

% The argument I showed you in your office shows that SL_2(Z_p) has only
% one possible simple quotient, namely SL_2(F_p)/{+I, -I}. Similarly for
% SL_2(Z_q).  Both of those quotients would be isomorphic to G, which is
% impossible.  Hence G can't be nonabelian.

% This proves that G is abelian. But the simple quotient of SL_2(Z_p) is
% nonabelian.  Hence the image of SL_2(Z_p) in G is trivial.  That means
% that K is contained in Q(mu_{p^{infty}}).  Similarly, K is contained
% in Q(mu_{q^{infty}}).  Ramification considerations give a
% contradiction.  Hence no such nontrivial K exists.
\end{remark}

% \begin{proposition}
% Suppose $Q\in E(K)$ has infinite order and that $n$
% is a positive integer divisible only by primes $\ell \in B(E)$.
% Then there exist a positive density of primes $p$ such that $p$
% is inert in $K$, splits completely in $K(E[n])/K$, 
% and the natural image of $Q$ in the quotient 
% group $E(\FF_{p^2})/(p+1)E(\FF_{p^2})$ has 
% order divisible by~$n$.
% \end{proposition}
% \begin{proof}
% [[sketch for now]]
% For each prime $\ell\mid n$, consider the extension
% $\QQ(E[n\ell])/\QQ(E[n])$.  This extension has Galois group isomorphic
% to the additive group $M_2(\FF_{\ell})$ of $2\times 2$ matrices
% with entries in $\FF_{\ell}$, equipped with its natural
% action on the quotient group $E[n\ell]/E[n] \isom \FF_{\ell}^2$. 
% Choose a matrix $A_{\ell} \in M_2(\FF_{\ell})$ with 
% characteristic polynomial that is irreducible over $\FF_{\ell}$.
% Let $n' = n \cdot \prod_{\ell\mid n} \ell$. 
% Then let 
% $$
%  \sigma \in \Gal(\QQ(E[n'])/\QQ(E[n]))
% \subset \Gal(\QQ(E[n'])/\QQ)
% $$
% be the element that maps to each $A_{\ell}$ under the natural map.
% Because $\QQ(E[n'])$ is linearly disjoint
% from $K$ [[why?!?]], there is an element $\tau\in \Gal(K(E[n'])/\QQ)$ 
% that maps to $\sigma$ and to the
% nontrivial element of $\Gal(K/\QQ)$.  By the Chebotarev density
% theorem, there is a positive proportion of primes $p$ such that
% $\Frob_p = [\tau]$.  For these primes $p$, we have that $p$
% is inert in $K$, splits completely in $K(E[n])/K$
% \end{proof}

% \section{A Refinement of Kolyvagin's Conjectures}\label{sec:kolyrefine}
% In this section, we make a new conjecture about the invariants
% $m_{\ell}$ that refines Kolyvagin's Conjecture~\ref{conj:C}, and which
% is motivated by the generalized Gross-Zagier formula \eqref{eq:ggz}
% and the theorems we prove will in Section~\ref{sec:maxwp}.  We also deduce
% implications of our conjecture for the structure of $V_{k,\ell^n}^f$.

% [[redo this section -- just what is prop~\ref{prop:conjd2}.  either conj it or
% prove it as a conditional theorem.]]

% \begin{proposition}\label{prop:conjd2}
% Assume Conjectures~\ref{conj:bsd} and \ref{conj:refined}, and let
% $f = f_{\ell}$.  For $\ell\in B(E)$ and all $k\gg 0$, we have 
% $V_{k,\ell^n}^f = \vphi(\ell^{m_{\ell,f}} E(\QQ))$,
% where $\vphi$ is the composition
% $E(\QQ) \to \H^1(\QQ,E[\ell^{n+k}]) \to H^1(K,E[\ell^{\infty}])$.
% \end{proposition}
% \begin{proof}
% [[need to give proof.]]
% \end{proof}

% [[ It just occurred to me that using his structure theorem one gets
% that $f_{\ell}$ is the smallest so that $V^a$ is nonzero, and this $V^a$ is
% nonzero so $a \geq f$.  That $a = \rank(V) - 1$ implies that $\rank(V)
% =a+1 \geq f+1$.  But if $V$ is in the minus part, then it has rank 
% $\leq f$ by one of Koly's structure theorems (I think!).  Thus conjecture D
% is all I really need.]]

% \begin{proposition}\label{prop:conjd0}
% Assume Conjectures~\ref{conj:bsd} and \ref{conj:refined}.
% Then  for all $\ell\in B(E)$, 
% Conjecture~\ref{conj:D} is true with $v=0$ and $a=f$.
% \end{proposition}
% \begin{proof}
%   Kolyvagin mostly states this, though without proof, on page 259 of
%   \cite{kolyvagin:structure_of_selmer}. [[This requires a real proof
%   which we will give here.]]
% \end{proof}

% \begin{proposition}\label{prop:conjd1}
%   Assume Conjectures~\ref{conj:bsd} and \ref{conj:refined}.  
% Let
%   $V$ be the group in Conjecture~\ref{conj:D} with $v=0$ and $a=f$.
%   Then for $\ell\in B(E)$, we have
% $$
%  V\tensor\ZZ_{\ell} = \ell^{m_f(\ell)} \cdot (E(\QQ)\tensor\ZZ_{\ell}).
% $$
% \end{proposition}
% \begin{proof}
%   Kolyvagin states this without proof on page 258 of
%   \cite{kolyvagin:structure_of_selmer}.
% [[Give a real proof here.]]
% \end{proof}


\section{Maximal Index Subgroups $W_p$}\label{sec:maxwp}
As above, we assume that $E$ is an elliptic curve over $\QQ$ with
positive analytic rank and that $K=\QQ(\sqrt{D})$ is a quadratic
imaginary field that satisfies the Heegner hypothesis and the
minimality hypothesis that $r_{\an}(E/\QQ) > r_{\an}(E^D/\QQ) \leq 1.$

Recall that for each inert prime $p$ of $K$ we defined a subgroup $X_p \subset
E(\FF_p)/(p+1)$ in Equation~\eqref{eqn:Xp} of
Section~\ref{sec:kolysubgroups}.  This was a group got by reducing
Kolyvagin points associated to all primes $\ell$ modulo a choice of
prime over $p$.  In this section, for all $\ell \in B(E)$ we
conditionally compute, in terms of $m_{\ell,f}$, the $\ell$-primary part
$X_p(\ell)$ of this subgroup $X_p\subset E(\FF_p)/(p+1)$.  We relate
our refinement of Kolyvagin's conjectures to the generalized
Gross-Zagier formula \eqref{eq:ggz}.  We also conditionally compute
$X_p$ in terms of $c\cdot \prod c_q\cdot \sqrt{\#\Sha(E/K)}$ using
Theorem~\ref{thm:kolythm1}.  We apply our description of $X_p$ to
prove that, up to primes not in $B(E)$, the subgroups $W_p$ with
$[E(K):W_p]$ maximal are all Gross-Zagier subgroups of $E(K)$.

\begin{proposition}\label{prop:main1}
Conjecture \ref{conj:D}
implies that for every $\ell \in B(E)$,
$$ 
   X_{p}(\ell) = \frac{p+1}{\ell^{v_{\ell}}}\cdot  \pi_p(\ell^{m_{\ell,f}} E(\QQ)),
$$
where $v_{\ell} = \ord_{\ell}(p+1)$. 
\end{proposition}
\begin{proof}
Let  $\Phi$ be the composite homomorphism
$$
\left(E(K_{\lambda})/ \ell^{v_{\ell}} E(K_{\lambda})\right)^{\Gal(K_{\lambda}/K)}
    \hookrightarrow \H^1(K_{\lambda}, E[\ell^{v_{\ell}}])^{\Gal(K_{\lambda}/K)}
\isom 
  \H^1(K, E[\ell^{v_{\ell}}]),
$$
and let $\delta:E(K) \to \H^1(K,E[\ell^{v_{\ell}}])$.
We are assuming Conjecture~\ref{conj:D}, so we may apply
Proposition~\ref{prop:conjD} Part~\ref{prop:D:6} 
(taking into account Lemma~\ref{lem:conjd}),  to see that
for all $k$ sufficiently large we have
$\delta(\ell^{m_{\ell,f}} E(\QQ)) = V_{k,\ell^{v_{\ell}}}^f$.
Thus
$$
  \delta(\ell^{m_{\ell,f}} E(\QQ))
   = \langle \Phi([P_{\lambda}]) : \lambda \in \Lambda_{\ell^{v_\ell+k}}^f \rangle.
$$
Let $i:E(K) \to \left(E(K_{\lambda})/ \ell^{v_{\ell}} E(K_{\lambda})\right)^{\Gal(K_{\lambda}/K)}$.
For any $Q\in E(K)$ we have 
$\frac{p+1}{\ell^{v_\ell}}\cdot \pi_p(Q)  =  \vartheta(i(Q))$
where $\vartheta$ is as in Proposition~\ref{prop:piplambdaell}. 
Since $\delta = \Phi \circ i$ and $\Phi$ is injective, 
the group $X_{k,p}(\ell)$ generated by all $\vartheta([P_{\lambda}])$  is
equal to $\vartheta(i(\ell^{m_{\ell,f}} E(\QQ)))$. 
Since this is true for all sufficiently large $k$, 
the proposition follows for $X_p$.
\qed\end{proof}

Theorem~\ref{thm:refined_evidence} below generalizes
\cite[Thm.~E]{kolyvagin:structureofsha} to arbitrary rank.  To prove
it we first prove some lemmas and make a definition.

\begin{lemma}\label{lem:main1}
  Suppose $A$ is a nonzero finitely generated free abelian group and
  $\vphi:A \to \ZZ/d\ZZ$ is a surjective homomorphism.  For every
  nonzero integer $c$ we have $[A:\vphi^{-1}(\vphi(c A))] = \gcd(c,d)$.
\end{lemma}
\begin{proof}
  Let $B = \vphi^{-1}(\vphi(c A))$.  We have $\vphi(\ker(\vphi)) = 0
  \subset \vphi(cA)$, so $\ker(\vphi) \subset B$.  Since
  $\ker(\vphi)\subset B$, the isomorphism $A/\ker(\vphi) \isom
  \ZZ/d\ZZ$ induces an isomorphism $A/B \isom (\ZZ/d\ZZ)/\vphi(B)$.
  But $\vphi$ is surjective, so $\vphi(B) =
  \vphi(\vphi^{-1}(\vphi(cA))) = \vphi(cA) = c\vphi(A) = c(\ZZ/d\ZZ)$,
  so $A/B \isom (\ZZ/d\ZZ)/(c(\ZZ/d\ZZ)) \isom \ZZ/\gcd(d,c)\ZZ$.
\qed\end{proof}

Recall (see page~\pageref{defn:be}) that $B(E)$ is a set of primes
that have certain good properties for $E$.  Below, for any integer $n$ we
either let $n'=\ell^{\ord_{\ell}(n)}$ be the $\ell$-part of $n$ or the
maximal divisor of $n$ divisible only by primes in $B(E)$, depending
on whether we are considering the first or second part of the following lemma.
\begin{lemma}\label{lem:t}
Assume $E(\QQ)$ has positive rank and let $t$ be a positive integer. 
\begin{enumerate}
\item If $\ell \in B(E)$ is such that  
$X_{p}(\ell) = \frac{p+1}{\ell^{v_{\ell}}}\cdot  \pi_p(t E(\QQ))$
for all inert primes~$p$, then
$$
  \max\{\ord_{\ell}([E(K):W_p]):\text{ all inert $p$}\} = \ord_{\ell}(t).
$$
\item If for all $\ell \in B(E)$ we have
$X_{p}(\ell) = \frac{p+1}{\ell^{v_{\ell}}}\cdot  \pi_p(t E(\QQ))$ 
for all inert primes $p$, then 
$$
  \max\{[E(K):W_p]':\text{ all inert $p$}\} = t'.
$$
\end{enumerate}
\end{lemma}
\begin{proof}
Let $p$ be any inert prime, and recall  that $p$ is a prime of good reduction, 
since all bad primes split in $K$.  By Lemma~\ref{lem:cyclic},
the odd part of the image of $\pi_p:E(K)\to E(\FF_p)/(p+1)$ is a cyclic group $\ZZ/n\ZZ$
for some integer $n$. 
Since $\pi_p(E^D(\QQ)) = 0$ (see Remark~\ref{rem:kertr}), we have 
$$
  \pi_p(t E(\QQ))' = \pi_p( t E(\QQ) + t E^D(\QQ))' = \pi_p(t E(K))',
$$
so by Proposition~\ref{prop:imteq}, 
$W_p' = \pi_p^{-1}(X_p)' = \pi_p^{-1}(\pi_p( t E(K) )')$.
Thus Lemma~\ref{lem:main1} implies that $[E(K)':W_p']$ is $\gcd(t, n)'$.
This proves that set of indexes $[E(K)':W_p']$ all divide $t'$. 

We show the maximum equals $t'$ by proving that 
there is a positive density of primes $p$ such that 
the $n$ above is divisible by $t'$.
By hypothesis, there is a point $P\in E(\QQ)$ of infinite order.
By Proposition~\ref{prop:bigord}, there exists a positive
density of primes $p$ that are inert in $K$ such that 
$\pi_p(P) \in E(\FF_{p})/(p+1)$ has order divisible by $t'$.
For such $p$, the $n$ above is thus divisible by $t'$, so 
$\gcd(t,n)'=t'$, which completes the proof.
\qed\end{proof}

Let 
\begin{equation}\label{eqn:well}
\  w_{\ell} = \sup(\{\ord_{\ell}([E(K):W_p])\,  : \, \text{ all inert } p\}) \leq \infty.
\end{equation}

\begin{lemma}\label{lem:index}
  Suppose $\ell\in B(E)$, that Conjecture~\ref{conj:D} is true for
  $E$, and assume that $p$ is an inert prime such that
  $\ord_{\ell}([E(K):W_p])$ is maximal in the sense that it 
equals $w_{\ell}$. Then $m_{\ell,f} =
  \ord_{\ell}([E(K):W_p])$.
\end{lemma}
\begin{proof}
\noindent{}We are assuming 
that Conjecture~\ref{conj:D} is true, so Proposition~\ref{prop:main1} applies
and gives an explicit formula for $X_p(\ell)$.
Namely, we may take $t=m_{\ell,f}$ in Lemma~\ref{lem:t}.
Also, by Conjecture~\ref{conj:D} (and Proposition~\ref{prop:conjD}) we have $E(\QQ)$
has rank at least~$1$.  The lemma then follows from Lemma~\ref{lem:t}.
\qed\end{proof}


\begin{theorem}\label{thm:refined_evidence}
  Suppose $\ell\in B(E)$, that Conjectures~\ref{conj:bsd} and
  \ref{conj:D} are true for $E$, and that $p$ is an inert prime such
  that $w_{\ell} = \ord_{\ell}([E(K):W_p])$, where $w_{\ell}$ is
as in \eqref{eqn:well} above.  Then $W_p$ satisfies the
  generalized Gross-Zagier formula \eqref{eq:ggz} up to a rational
  factor that is coprime to $\ell$ if and only if
  Conjecture~\ref{conj:refined} is true for $\ell$.
\end{theorem} 
\begin{proof}
We are assuming Conjecture~\ref{conj:D}, which implies Conjecture~\ref{conj:A},
so we may apply Theorem~\ref{thm:kolythm1}, which has Conjecture~\ref{conj:A}
as a hypothesis. 
Let $b_k$ be as in Theorem~\ref{thm:kolythm1} for our
  given prime $\ell$.  Theorem~\ref{thm:kolythm1} implies that
\begin{align*}
 \#\Sha(E/K)(\ell) &= 
\#((\ZZ/b_f \ZZ)^{2}
     \oplus (\ZZ/b_{f+1}\ZZ)^{2}
     \oplus (\ZZ/b_{f+2}\ZZ)^{2}
\oplus
\cdots)\\
 &= (b_f \cdot b_{f+1} \cdots )^2\\
 & = \ell^{2(m_{\ell,f} - m_{\ell,f+1} + m_{\ell,f+1} - m_{\ell,f+2} + m_{\ell,f+2} - \cdots)} \cdots\\
 &= \ell^{2(m_{\ell,f} - m_{\ell})},
\end{align*}
so $ m_{\ell,f} - m_{\ell} = \ord_{\ell}(\sqrt{\#\Sha(E/K)(\ell)})$.

% When $E$ has analytic rank $1$, Corollary~\ref{cor:gzrank1} 
% and Conjecture~\ref{conj:bsd} imply that 
% \begin{align*}
% m_{\ell,f} &= \ord_{\ell}([E(K): \ZZ P_1])\\ 
%  &= \ord_{\ell}\left(c \prod c_q \cdot \sqrt{\#\Sha(E/K)(\ell)}\right)\\
%  &= \ord_{\ell}\left(c \prod c_q\right) + \ord_{\ell}(\sqrt{\#\Sha(E/K)(\ell)})\\
%  &= \ord_{\ell}\left(c \prod c_q\right) + m_{\ell,f} - m_{\ell}
% \end{align*}
% Subtracting $m_{\ell,f}$ from both sides implies that $m_{\ell} = \ord_{\ell}(c \prod c_q)$.  
%Now assume $E$ has analytic rank $\geq 1$ and 


We will now show that the generalized Gross-Zagier formula \eqref{eq:ggz}
holds up to a rational factor that is coprime to $\ell$ if and only if
Conjecture~\ref{conj:refined} that $m_{\ell} = \ord_{\ell}(c \prod
c_q)$ is true for $\ell$.   We will repeatedly use 
Lemma~\ref{lem:index} that $m_{\ell,f} = \ord_{\ell}([E(K):W_p])$.

First, suppose that the generalized Gross-Zagier formula
\eqref{eq:ggz} holds up to a rational factor that is coprime
to~$\ell$.  Proposition~\ref{prop:ggz} combined with
Conjecture~\ref{conj:bsd} (that $\Sha_{\an}=\#\Sha$), 
implies that 
this hypothesis means that $\ord_{\ell}([E(K): W_p]) = \ord_{\ell}\left(c \prod c_q \cdot
  \sqrt{\#\Sha(E/K)(\ell)}\right)$.  Thus:
\begin{align*}
m_{\ell,f} &= \ord_{\ell}([E(K): W_p])\\ 
 &= \ord_{\ell}\left(c \prod c_q \cdot \sqrt{\#\Sha(E/K)(\ell)}\right)\\
 &= \ord_{\ell}\left(c \prod c_q\right) + \ord_{\ell}(\sqrt{\#\Sha(E/K)(\ell)})\\
 &= \ord_{\ell}\left(c \prod c_q\right) + m_{\ell,f} - m_{\ell},
\end{align*}
where in the last equality we use the formula for $\#\Sha(E/K)(\ell)$ that we derived
above using Theorem~\ref{thm:kolythm1}.
Subtracting $m_{\ell,f}$ from both sides shows that $m_{\ell} = \ord_{\ell}(c \prod c_q)$.

Conversely, suppose that $m_{\ell} = \ord_{\ell}(c \prod c_q)$.
From Theorem~\ref{thm:kolythm1} we have
$$
  m_{\ell,f} - m_{\ell} = \ord_{\ell}(\sqrt{\#\Sha(E/K)(\ell)}),
$$
so
$$
\ord_{\ell}([E(K):W_p]) = m_{\ell,f} = \ord_{\ell}\left(c \prod c_q\right) + m_{\ell,f} - m_{\ell}
 = \ord_{\ell}\left(c \prod c_q \cdot \sqrt{\#\Sha(E/K)(\ell)}\right).
$$
Proposition~\ref{prop:ggz} then implies that $W_p$ satisfies the 
generalized Gross-Zagier formula up to a rational factor coprime to $\ell$.
\qed \end{proof}


For any integer $n$, let $n'$ denote the maximal divisor of $n$ that
is divisible only by primes in $B(E)$, and for any abelian group
$A$, let $A' = A\otimes\ZZ[1/b]$, where $b$ is the product of the
finitely many primes not in $B(E)$.
Let $$T = c  \cdot \prod_{q \mid N} c_q \cdot \sqrt{\#\Sha(E/K)}.$$

\begin{proposition}\label{prop:imteq}
Conjectures \ref{conj:D} and \ref{conj:refined}
together imply that $X_p' = \pi_p(T E(\QQ))'$.
\end{proposition}
\begin{proof}
Using the calculation in the first paragraph of the proof of Theorem~\ref{thm:refined_evidence}
along with Conjecture~\ref{conj:refined} combined with
Theorem~\ref{thm:kolythm1}, shows that for every  $\ell\in
B(E)$, we have
$$
  m_{\ell,f} = \ord_{\ell}(T).
$$
Since the integers $T/\ell^{m_{\ell,f}}$ and $(p+1)/\ell^{v_\ell}$,
for $v_{\ell}=\ord_{\ell}(p+1)$, both act as automorphisms on any
$\ell$-primary group, 
\begin{align*}
  \pi_p(T E(\QQ))(\ell) &= \left(\frac{T}{\ell^{m_{\ell,f}}} \cdot \pi_p(\ell^{m_{\ell,f}} E(\QQ))\right)(\ell)\\
   &= \pi_p(\ell^{m_{\ell,f}} E(\QQ))(\ell) \\
   &= \frac{p+1}{\ell^{v_{\ell}}}\cdot  \pi_p(\ell^{m_{\ell,f}} E(\QQ)) = X_p(\ell),
 \end{align*}
 where the last equality uses Proposition~\ref{prop:main1} (which
 assumes that Conjecture~\ref{conj:D} is true).  We conclude that
 $X_p' = \pi( T E(\QQ))'$.
\end{proof}



Theorem~\ref{thm} is a partial converse to Theorem~\ref{thm:refined_evidence}.

\begin{theorem}\label{thm}
  Assume that $E(\QQ)$ has positive rank.  Then
  Conjectures~\ref{conj:D} and \ref{conj:refined} together imply that the
  maximum index $[E(K)':W_p']$ over all inert $p$ is $(c \cdot \prod
  c_q \cdot \sqrt{\#\Sha(E/K)})'$.
\end{theorem}   
\begin{proof}
The conjectures we're assuming allow us to use Proposition~\ref{prop:imteq}
and hence take $t=T$ in Lemma~\ref{lem:t}.  This proves the theorem. 
\qed\end{proof}

\noindent{\bf Conclusion:}
By Proposition~\ref{prop:ggz}, if $W_p'$
has maximal index in $E(K)'$, then  imply that we have an equality
$$
\frac{L^{(r)}(E,1)}{r!} = \frac{\|\omega\|^2}{c \cdot \sqrt{|D|}} \cdot \Reg(W_p),
$$
up to powers of primes not in $B(E)$.
Thus the $W_p'$ of maximal index satisfy this generalized Gross-Zagier formula.

\begin{conjecture}
  If $W\subset E(K)$ is {\em any} Gross-Zagier subgroup of index
  $\ell^{w_{\ell}}$, then there exists an inert prime $p$ such that
  $W_p'$ equals $W'$.
\end{conjecture}




\section{Existence of Gross-Zagier Subgroups}\label{sec:gzexist}
Let $E$, $K$, etc., be as in Section~\ref{notation}, and let
$$ 
  t = c \cdot \prod_{q\mid N} c_q \cdot \sqrt{\#\Sha(E/K)_{\an}}.
$$

In this section we investigate the analogue of the conjectures on
pages 311--312 of \cite{gross-zagier}.  In particular, the existence
of any Gross-Zagier subgroup for $E(K)$ combined with the BSD
conjecture implies that $\#E(K)_{\tor} \mid t$.  The main theorem of
\cite{gross-zagier} thus led Gross-Zagier to make the following
conjecture.
\begin{conjecture}[Gross-Zagier]\label{conj:gz}
  If $E(K)$ has rank $1$, then the integer $t$ is divisible by $\#E(\QQ)_{\tor}$.
\end{conjecture}

\begin{proposition}\label{prop:anyggz}
Assume the BSD formula.  If there exists any subgroup $W$ of $E(K)$ such that 
the generalized Gross-Zagier formula \eqref{eq:ggz} holds for $W$, then
$\#E(K)_{\tor}\mid t$. Note that we do not assume $W$ is torsion free.
\end{proposition}
\begin{proof}
Let $W$ be such a subgroup.
Arguing as in the proof of Proposition~\ref{prop:ggz}, we see that
$$
  \#E(K)_{\tor}^2 \cdot (\Reg(W)/\Reg(E/K))
 = c^2 \cdot \left(\prod_{q\mid
    N}c_q\right)^2 \cdot \Sha_{\an} = t^2.
$$ 
The quotient $\Reg(W)/\Reg(E/K)$ is a square integer, so
taking square roots of both sides yields the claim.

\qed\end{proof}

Because of Proposition~\ref{prop:anyggz}, we view the divisibility
$\#E(K)_{\tor} \mid t$ as a sort of {\em ``litmus test''} for whether
there could be a generalization of the Gross-Zagier formula in
general.  First, we observe that the most naive generalization of
Conjecture~\ref{conj:gz} to higher rank is {\em false} (!), as the
following example shows.
\begin{example}\label{ex:gzcounter}
  Let $E$ be the curve 65a of rank $1$ over $\QQ$ given by $y^2 + xy =
  x^3 - x$ and let $D=-56$. Then $\#\Sha_{\an}(E/K)=\prod c_q = c = 1$, so
  $t=1$, but $\#E(\QQ)_{\tor} = 2$.  Here $E^D(\QQ)$ has rank $2$, so
  $\rank(E(K))=3$, and the rank hypothesis of Conjecture~\ref{conj:gz}
  is not satisfied.
\end{example}

\begin{proposition}\label{prop:gzcriterion}
Suppose $\rank(E(\QQ))> 0$ and that $t$ is a positive integer. 
Then there exists a Gross-Zagier subgroup $W\subset E(K)$ if and only if
$\#E(K)_{\tor} \mid t$.
\end{proposition}
\begin{proof}
Suppose $W\subset E(K)$ is a Gross-Zagier subgroup.  Then $[E(K):W] = t$. 
By hypothesis $W$ is torsion free, so $E(K)_{\tor} \hookrightarrow E(K)/W$, so 
$\#E(K)_{\tor} \mid \#(E(K)/W) = t$.  

Conversely, suppose that
$\# E(K)_{\tor} \mid t$, and note that by hypothesis $E(\QQ)$ has positive rank. 
The group $E(K)/(E^D(\QQ) + E(K)_{\tor})$ is thus a finitely generated
infinite abelian group, so has subgroups of all index.  In particular,
it has a subgroup $W'$ such that the quotient
by $W'$ is cyclic of order $t/\#E(K)_{\tor}$.
Let $\tilde{W}$ be the inverse image of $W'$ in $E(K)$, so 
$E(K)_{\tor}, E^D(\QQ) \subset \tilde{W}$,
and $[E(K):\tilde{W}] = t/\#E(K)_{\tor}$. 
Since $\tilde{W}$ is finitely generated, there exists a torsion free
subgroup $W\subset \tilde{W}$ such that $W\oplus E(K)_{\tor} = \tilde{W}$. 
Then $$
[E(K):W] = \#E(K)_{\tor} \cdot [E(K) : \tilde{W}]
          = \#E(K)_{\tor} \cdot \frac{t}{\#E(K)_{\tor}} = t.
$$
\qed\end{proof}

Elsewehere in this paper, for technical reasons in order to apply
Kolyvagin's theorems, we made a minimality hypothesis on
$r_{\an}(E^D/\QQ)$, and based on extensive numerical data, we
conjecture that this is the right hypothesis to guarantee the existence
of Gross-Zagier subgroups $W\subset E(K)$.
\begin{conjecture}\label{conj:gzgen}
  If $r_{\an}(E/\QQ) > r_{\an}(E^D/\QQ) \leq 1$, then $\#E(K)_{\tor}
  \mid t$.  In particular, there exists a Gross-Zagier subgroup
  $W\subset E(K)$.
\end{conjecture}

We obtain evidence for Conjecture~\ref{conj:gzgen} using Sage\footnote{Running
on hardware purchased using National Science Foundation Grant No. DMS-0821725.}
\cite{sage, mwrank, pari}, Cremona's tables
\cite{cremona:onlinetables}, Proposition~\ref{prop:gzcriterion}, and
assuming the Birch and Swinnerton-Dyer conjecture.  More precisely, we
check that Conjecture~\ref{conj:gzgen} is ``probably true'' for every
elliptic curve of rank $\geq 2$ and conductor $\leq 130,000$ and the
first three $D$ that satisfy the Heegner hypothesis, except possibly
for the triples $(E,D,\#E(K)_{\tor})$ in Table~\ref{table:triples}
where the computation of the conjectural order of $\#\Sha(E/K)$ took
too long.
\begin{table}[ht]
\caption{All triples up to conductor 130,000 where we did not yet verify Conjecture~\ref{conj:gzgen}\label{table:triples}}
\begin{tabular}{|c|}\hline
$ (8320e1,-191,2),  (9842d1,-223,3),  (9842d1,-255,3),  (9842d1,-447,3),
(74655j1,-251,3),  $\\

$
(87680a1,-119,2),  (87680a1,-151,2),  (87680b1,-119,2), 
(87680b1,-151,2),  (89465a1,-51,2),  $\\

$
(89465a1,-59,2),  (89465a1,-71,2),  (95545b1,-191,2),  (95545b1,-219,2),
(104585b1,-139,2),  $\\

$
(104585b1,-179,2),  (104585b1,-191,2),  (114260a1,-231,2), 
(114260a1,-239,2),  (114260a1,-431,2),  $\\

$
(122486a1,-103,3),  (122486a1,-55,3),  (122486a1,-87,3), 
(126672r1,-335,2),  (126672r1,-647,2),  $\\

$
(126672r1,-719,2),  (129940a1,-111,2),  (129940a1,-71,2), 
(129940a1,-79,2)  $\\
\hline
\end{tabular}
\end{table}

In our computations, we considered the first three Heegner $D$, {\em
  without} making the condition $r_{\an}(E^D/\QQ) \leq 1$.  The
conjecture is {\em false} without the hypothesis that
$r_{\an}(E^D/\QQ) \leq 1$, as Example~\ref{ex:gzcounter} above shows.
Moreover, we found two further similar examples in which, however, $E$
has rank $2$ and $E^D$ has rank $3$.  First, for the curve $E$ with
Cremona label 20672m1, equation $y^2 = x^3 - 431x - 3444$ and
$D=-127$, we have $\rank(E(\QQ))=2$, $\rank(E^D(\QQ)) = 3$, and
$\#E(K)_{\tor} = 2$, but $t = 1$.  A second example is $E$ given by
18560c1 and $D=-151$, in which again $\rank(E(\QQ))=2$,
$\rank(E^D(\QQ)) = 3$, $\#E(K)_{\tor} = 2$, but $t=1$.


This was a large computation that relies on a range of nontrivial
computer code, which we carried out as follows.  First we computed
$\#E(K)_{\tor}$ for each of the 78,420 elliptic curve of conductor
$\leq$130,000 with rank $\geq 2$ and the first three Heegner $D$. We
then determined whether $\#E(K)_{\tor}$ divides $c \cdot \prod c_q$.
Since we are verifying that something divides $c \cdot \prod c_q$,
there is no loss at all in assuming Manin's conjecture that $c=1$ for
the optimal quotient of $X_0(N)$.  We then computed the Manin constant
$c$ for non-optimal curves by finding a shortest isogeny path from the
optimal curve in the isogeny graph of $E$ (there is unfortunately a
small possibility of error in computation of the isogeny graph, due to
numerical precision used in the implementation).  We found only $37$
remaining curves $E$ of rank $\geq 2$ such that $\#E(K)_{\tor} \nmid
c\cdot \prod c_q$, and $37\cdot 3=111$ corresponding pairs $(E,D)$.
It turns out that all of these curves are optimal hence have $c=1$.
For each of these pairs $(E,D)$ we attempted to compute
$\#\Sha(E/K)_{\an}$ using Conjecture~\ref{conj:bsd} and some results
of \cite{bsdalg1}, and the computation finished in all but $29$ cases.
The main difficulty was computing $\Reg(E/K)$ in terms of
$\Reg(E/\QQ)$ and $\Reg(E^D/\QQ)$ by saturating the sum of $E(\QQ)$
and $E^D(\QQ)$ in $E(K)$.  Computing $E^D(\QQ)$ was sometimes very
difficult, since $E^D$ has huge conductor and rank $1$, and this
sometimes took as long as a day when it completed.  For more details,
the reader is urged to read the source code of the Sage command
\verb|heegner_sha_an| in Sage-3.4.1 and later.

%[[TODO: Add table of the examples with forced nontrivial Sha.]]

\bibliography{biblio}
  
\end{document}


